% =====================================================================
%  R. Nielsen, DR. H. MARTENSENS DOGMATISKE OPLYSNINGER BELYSTE.
%  Kjøbenhavn: C. A. Reitzel (Bianco Luno), 1850.  75 pp.  Fraktur.
%
%  TRANSCRIPTION -- GENERATED; DO NOT EDIT BY HAND.  Rebuild, chained (&&), from the repo root:
%    python3 ebook-method/martensens-oplysninger/draft.py --emit --force
%      this template + the scan's text layer (the Royal Library's ABBYY OCR in KB's scan of copy
%      DA 1.-2.S 3 8°; ~/bibliotek/Nielsen, Rasmus/1850-martensens-oplysninger.pdf, see SCANS.tsv);
%    python3 ebook-method/collate/check.py martensens-oplysninger --witness ocr --apply
%      the decisions in decisions.tsv, taken against tesseract's independent reading of the images;
%    python3 ebook-method/martensens-oplysninger/oepass.py --apply
%      the ø pass and the word fixes of the unattested-word sweep (word-fixes.tsv);
%    python3 ebook-method/martensens-oplysninger/finish.py --apply
%      the footnotes, corrections across markup, Latin (\textit), letterspacing (\emph), PRINTED AS IS marks.
%
%  WORD-ACCURACY (2026-10-06).  No e-book and no earlier transcription exist.  The draft (KB's
%  OCR layer) was collated word by word, with punctuation, quote kinds and word division, against
%  tesseract (tessdata_best script/Fraktur): 2985 disagreements.  1593 were settled by rule
%  (rules.py, each decision named): each OCR's known Fraktur glyph faults undone, accepted only
%  where the cheapest agreement is a word attested in at least two other books of the corpus, or
%  where one OCR's literal reading is the other's after at most two of the other's known faults;
%  a third reading (tesseract deu_latf) voting under the witness's word box.  The rest were read
%  at the image: 851 crops read BLIND by six readers (pilot + 5, controls 12 of 12), decided by
%  script (readsbox.py, reads.py) or by hand from the readers' transcriptions; a second round of
%  97 crops for what remained.  7 rule decisions found wrong afterwards (tesseract's x is sometimes
%  the print's: Paradox, strax, Excerpt) were corrected and the rule narrowed.
%  Ø: all three OCRs read the print's ø as o (the draft had 7 ø in 73 pp.).  oepass.py restores
%  it where the o-form is no word (349 tokens) and, for pairs that are both words (for/før,
%  bor/bør, tor/tør, fore/føre ...), from 229 tight crops of the word alone, without context
%  (controls 4 of 4): 33 ø.  The print has „før“ twice where sense needs „for“ (p. 73), checked
%  at the image: PRINTED AS IS.
%  Then at the image: the unattested-word sweep (130 words found in no other book of the corpus,
%  128 read: 78 OCR faults mended, 5 misprints kept and marked); every footnote read whole
%  (eleven pages); the title page and motto page; the Latin and Greek; letterspacing (245
%  candidate words measured on the images, 4 spaced, and 9 more from the readers' notes and the
%  main session's own look).  Spot-read in full, pp. 37 and 62: 3 faults found, 2 mended (a word
%  split by a lost line-end hyphen; „sor“ for „for“ over a page turn).
%  The translators' reading: 134 odd spots reported; 59 crops read at the image (T01-T59,
%  .parts/reads/ts-reads.tsv): about 30 faults mended (ø, s/f, a doubled word, two false paragraph
%  breaks, a stray „15“ signature mark), 6 more misprints kept and marked; „hævder“ (p. 62)
%  confirmed as printed.  Dropped text: check.py --dropped (witness words the transcription
%  lacks, two or more in a run) found a whole line the layer had lost (p. 4, „efter at have
%  slaaet betydeligt“) and three shorter runs (pp. 8, 20, 56); all four mended from the image.
%  NOT ESTABLISHED: misprints that both OCRs read alike as real words, beyond the two spot pages
%  (no full image pass); letterspacing outside the candidates read (short words; the extent of
%  the spaced sentences on pp. 47-48 rests on single lines seen by the main session); in the
%  footnotes, where the note reader and the collation differed, the collation's word was kept
%  (d/b and o/ø are hard to see in the small type).
%  PAGE MARKERS: 73 of 73 exact (check.py --pages; p. 51 turns inside „ideologisk“, split with %).
%
%  CONVENTIONS (as Om „Den gode Villie“).  Body in Fraktur; aa (not å); ø as printed.  Emphasis
%  is letterspacing -> \emph{}.  Latin set in antiqua -> \textit{}.  Greek in Greek type
%  (textalpha).  Fraktur I/I are one sort: I for I-words (Idee, Ikke), I for J-words (Jeg).
%  The double hyphen at a line end is resolved; inside a compound it is "-".  The reversed c
%  „ɔ:“ (id est) as printed.  "% --- p. N ---" + \opage{N} at the first word of printed page N.
%  Footnotes *), **) per page, as printed.  Folios, signature marks (p. 49 "4", p. 51 "4*") and
%  the library's stamps and pencil marks are not transcribed.  Print governs; misprints are kept
%  and marked PRINTED AS IS.
% =====================================================================
\documentclass[12pt,a4paper,openany]{book}
\usepackage[utf8]{inputenc}
\usepackage[T1]{fontenc}
\usepackage{libertinus}
\usepackage{libertinust1math}
\usepackage{textalpha}
\usepackage{graphicx}
\DeclareUnicodeCharacter{0254}{\reflectbox{c}}  % ɔ (U+0254) = reversed c
\usepackage[danish]{babel}
\usepackage[protrusion=true,expansion=true]{microtype}
\usepackage{setspace}
\usepackage[margin=1.3in]{geometry}
\usepackage{fancyhdr}
\usepackage{hyperref}
\hypersetup{pdftitle={Dr. H. Martensens dogmatiske Oplysninger belyste}, pdfauthor={Rasmus Nielsen}, hidelinks}
\setlength{\headheight}{15pt}
\pagestyle{fancy}
\fancyhf{}
\fancyhead[LE,RO]{\thepage}
\fancyhead[C]{\textit{R. Nielsen: Dr. H. Martensens dogmatiske Oplysninger belyste (1850)}}
\renewcommand{\thefootnote}{\ifcase\value{footnote}\or *)\or **)\or ***)\else
  \arabic{footnote})\fi}
\setstretch{1.3}
\usepackage{marginnote}
\newcommand{\opage}[1]{\marginnote{\footnotesize\textit{[#1]}}}
\newcommand{\apage}[1]{\marginnote{\footnotesize\textit{[#1]}}}

\begin{document}

% --- front matter, read at the image 2026-10-06 (.parts/reads/final-reads.txt, JOB A).
% p. [1], title.  "Dr." is in antiqua; "Martensens", "Reitzel." and "Bianco Luno." are letterspaced.
% The ornamental rule (a central flower between tapering rules) is drawn as a plain rule; KB's stamp is not transcribed.
\begin{titlepage}
\begin{center}
\vspace*{3em}
{\large Dr. H. \emph{Martensens}}\\[1em]
{\Huge Dogmatiske Oplysninger}\\[1.5em]
belyste\\[0.8em]
af\\[0.8em]
Prof. R. Nielsen.\\[2em]
\rule{8em}{0.4pt}\\[2em]
{\large Kjøbenhavn.}\\[0.5em]
Hos Universitetsboghandler C. A. \emph{Reitzel.}\\[0.5em]
{\small Trykt hos Kgl. Hofbogtrykker \emph{Bianco Luno.}}\\[0.5em]
1850.
\end{center}
\end{titlepage}
% p. [2], the motto, in small type at mid-page.
\thispagestyle{empty}
\vspace*{12em}
\begin{center}
„I dit Lys skulle vi see Lys!“\\[0.4em]
\hspace*{6em}{\small En dogmatisk Oplysning.}
\end{center}
\clearpage


% ===== TEXT (pp. 3--75) =====
% --- p. 3 ---
\setcounter{footnote}{0}%
\opage{3}Min „undersøgende Anmeldelse“ har altsaa gjort sin Virkning;
Dr. Martensen, der i Forordet til sin Dogmatik lod
til at ville affærdige visse Punkter paa en temmelig tvetydig
og flygtig Maade, har altsaa dog omsider indseet Nødvendigheden
af at erklære sig med lidt større Bestemthed,
og har da nu ogsaa i sine „dogmatiske Oplysninger“ oplyst
os Adskilligt, som uden den „undersøgende Anmeldelse“ vistnok
var forblevet uoplyst.

At disse med „saamegen Klarhed“ affattede Oplysninger
maae have vakt Sensation i den theologiske Læseverden
og aabnet Manges Øine, tør jeg saameget mindre
betvivle, som ogsaa jeg afnødes den Bekiendelse, at „Klarheden“
i disse dogmatiske Oplysninger langt overgaaer, hvad
jeg i saa Henseende havde ventet. Kan altsaa Dr. M. for
sit Vedkommende være tilfreds med sine Oplysninger, da er
jeg efter Omstændighederne i Sandhed ogsaa tilfreds; ja
denne min Tilfredshed har netop bevirket, at en forholdsviis
lang Tid allerede er hengaaet, inden jeg igjen kunde
beqvemme mig til at afbryde min Ro, og tage fat paa
Sagen.

 Hvad der da især maa fremhæves som umiskjendeligt
% --- p. 4 ---
\setcounter{footnote}{0}%
\opage{4}Fortrin ved de „dogmatiske Oplysninger“, er den uindtagelige
Position, hvori Dr. M. her stiller sig til „Kritiken“, og den
øiensynlige Overlegenhed, hvormed han hævder sit dogmatiske
ποῦ στῶ ligeoverfor sine afmægtige „Recensenter.“

Her kan da naturligviis ikke være Tanke om, at Dr.
M. skulde være bragt i nogensomhelst Forlegenhed eller have
endog en eneste Indrømmelse at gjøre sine Modstandere.
Tvertimod, med hans dogmatiske Speculation er Alt i sin
bedste Orden: Dogmatiken, dens Udførelse, dens Indledning,
dens Forord, Alt er her saa sikkert, saa begrundet, saa sammenhængende,
at Den, der nu ikke kan see, at „Kritikens“
Indsigelser ere ligesaa ubeføiede, som Dr. Ms Recensenter
incompetente, maa mangle al „Evne til umiddelbar og levende
Ideeanskuelse.“

Man har nu rigtignok meent, at en Kritik over Dr.
Ms videnskabelige Behandling af de christelige Dogmer virkelig
var kommen i Gang; men, om det nu ligger i Mangel
paa „Evne til umiddelbar og levende Ideeanskuelse“,
ifølge de dogmatiske Oplysninger, er det en reen Indbildning;
Kritiken er endnu ikke begyndt. „Jeg kan“ --- siger
Dr. M. (S. 7) --- kun tillægge hele denne Critik en blot
foreløbig Charakteer, ja maa finde, at Critiken over min
Dogmatik, saafremt en saadan overhovedet turde forventes,
endnu ikke kan siges at være begyndt.“\dots{}

Man har endvidere meent, at der i Mag. Kierkegaards
mange Skrifter, og navnlig i Pseudonymen „Johannes Climacus“,
var Adskilligt, som vel turde komme i Betragtning,
naar det gjaldt om at vurdere en Dogmatiker, som,
efter at have slaaet betydeligt af i Speculationen, dog ikke
har kunnet beqvemme sig til at opgive Speculationen; men ---
% --- p. 5 ---
\setcounter{footnote}{0}%
\opage{5}det maa vist ligge i Mangel paa „Evne til umiddelbar og
levende Ideeanskuelse“ --- ogsaa dette er, ifølge de dogmatiske
Oplysninger, en reen Vildfarelse. Dr. M. har (S.
13) aldeles ingen Brug for de „Læresætninger“, man som
Uddrag af Pseudonymen har villet meddele ham, efterdi han
paa disse „maa anvende det gamle Ord, at det Sande deri
ikke er Nyt, og det Nye ikke Sandt.“

Man har endelig meent, at den evangeliske Opfattelse
af Troen, som det for den naturlige Forstand uigjennemtrængelige
Paradox, maatte, efter Alt hvad den nyere Tid har
oplyst om Ueensartetheden af Tro og Viden, naar Begrebet
skjærpedes, dog sætte den speculerende Dogmatiker i nogen
Forlegenhed; men --- ogsaa dette maa vist hidrøre fra Mangel
paa „Evne til umiddelbar og levende Ideeanskuelse ---
ogsaa dette er, ifølge de dogmatiske Oplysninger, en reen
Vildfarelse. Knuden stikker i, at man har forsømt at oversætte
Ordet Paradox i Modersmaalet. Dr. M., der „i
Almindelighed pleier, naar en vanskelig Skoleterminologie
foreligger ham, at opsøge det tilsvarende Udtryk i Modersmaalet“,
har da ogsaa, ved at give adskillige ret passende
Oversættelser af dette vanskelige Ord, med sædvanlig Grundighed
udfundet, at „Paradoxet“, naar man blot søger de
tilsvarende Udtryk i Modersmaalet, aldeles Intet udsiger,
som jo meget vel lader sig forene med --- „den dogmatiske
Speculation.“

Ere nu Recensenternes kritiske Indvendinger virkelig
saa intetsigende, som sees kan i de dogmatiske Oplysninger,
saa kan man da heller ikke være uenig om, hvorledes de
Personer, der have kunnet fremsætte slige Indvendinger,
maae blive at behandle. Ogsaa dette har Dr. M. forstaaet.
% --- p. 6 ---
\setcounter{footnote}{0}%
\opage{6}Istedetfor de sædvanlige Omskrivninger, hvorved man ellers
i den litterære Polemik søger at blotte sin Modstander uden
dog at krænke ham moralsk, betjener Dr. M. sig her med
Alvor og Myndighed af den „ligefremme Meddelelse“, og
lader uden Omsvøb det ærede Publicum vide, hvad hans
Modstandere ere for Folk.

Den Ene af disse Modstandere, nemlig Mag. Stilling,
er da saaledes, ifølge de dogmatiske Oplysninger (Jfr. S.
7 -- 8), et i speculativ-dogmatisk Forstand aldeles forulykket
Subject, et Menneske, der i den Grad mangler „positive
Kundskaber og et paa religiøs Livserfaring grundet Bekjendtskab
med Aabenbaringens Sandheder, at han endog
har meent paa et Par Blade at kunne affærdige Treenighedslæren,
Incarnationslæren, Ubiqvitetslæren, Læren om
Syndefaldet og Arvesynden, om Frihed og Naade, foruden
en heel Cyclus af andre Dogmer“; en Dilettant i Logiken,
der i den Grad er berøvet „Evne til umiddelbar og levende
Jdeeanskuelse“, at Dr. M. slet ikke kan indlade sig med ham,
især da han --- hvad ikke kan Andet end støde Dogmatikerens
æsthetiske Sands, har --- „uvadskede Hænder“ og ligner
„Erasmus Montanus!“ Med denne „dogmatiske Oplysning“
er Mag. St. altsaa een Gang for alle afviist; Dr. M. veed ikke,
om Andre ville „tage nærmere Hensyn til Hr. Magisterens
Skriveri“, Hs. Høiærværdighed har (Oplysn. S. 8) for sit
Vedkommende strax lagt det hen, for ikke mere at optage det.“

Den Recensent, der med et Apparat af Beviser og Grunde,
som det, hvormed Mag. St. har udstyret ovennævnte
Kritik, kan opnaae, at den kritiserede Forfatter behandler
ham saa forekommende, som Dr. M. i nærværende Tilfælde
behandler Mag. Stilling, kan imidlertid gratulere sig; thi det
% --- p. 7 ---
\setcounter{footnote}{0}%
\opage{7}er jo umiskjendeligt Beviis paa, at Kritiken har virket tilgavns.
Mag. St. har saaledes denne Gang unegtelig været heldigere
end jeg; hans Kritik har aabenbart gjort en grundigere
Virkning, end min, hvorfor Dr. M. da ogsaa behandler
ham med øiensynlig Forkjærlighed. Alligevel er jeg dog ikke
misundelig; jeg under gjerne Mag. Still. al den litterære
Anerkjendelse, han med Rette fortjener, og har desuden al
Grund til efter Omstændighederne at være tilfreds med det
moralske Skudsmaal, de „dogmatiske Oplysninger“ tilkjende
mig; her er det, det lyder saaledes: „Salige ere de, der
ikke ville paanøde Andre deres Viisdom!“ Dette Ord af
„den berømte Dogmatiker, den gamle Daub“, maatte Dr.
Martensen (S. 9) uvilkaarligt mindes, da han saae „den
Fremfusenhed“, hvormed Prof. N. vil paanøde ham den
Viisdom, han (Prof. N.) mener at have fundet i Mag.
Kierkegaards Skrifter. Prof. N., der dog „selv maa erkjende,
at han ved et Lykketræf er kommen over paa det
nye Standpunkt“, er nu tillige saa fortrydelig over, at Dr.
M. ikke ogsaa har skiftet Standpunkt, at han endog „med
en til Fanatisme grændsende Heftighed“ vil give Hs. Høiærværdighed
„Syn paa Problemet.“ Den fremfusende, paatrængende,
fanatiske Prof. N. er dog endnu (Dogm. Oplysn.
S. 11) ikke selv bleven hjemme paa sit nye Standpunkt, og
Dr. M. kan derfor (S. 11) ikke Andet end finde, at det er
gaaet ham, som det pleier at gaae Begyndere, der antage
en ny philosophisk Lære, at de ikke blot tænke med deres
Mesters Tanker, men ogsaa tale i disses Ord, og for en Tid tabe
Brugen af deres egen Tunge. Den fremfusende, paatrængende,
fanatiske Prof. N. er endnu selv (Dogm. Oplysn. S. 11)
„saaledes indviklet i Svøbet af den nye Viisdom“, og „be%
% --- p. 8 ---
\setcounter{footnote}{0}%
\opage{8}væger sig endnu saa besværlig, saa discipulært i denne Viisdoms
første Elementer, at han ikke engang formaaer at foredrage
Dr. Martensen sine Læresætninger med sine egne
Ord.“ --- „At Prof. N. ikke var synderlig stærk i at anerkjende
specifiske Forskjelle og iagttage Grændser“, har Dr. M.
(Opl. S. 62) længe vidst; men nu gaaer Forvirringen i dette
fremfusende, fanatiske Hoved saavidt, at Prof. N. ikke længere
gjør Forskjel imellem Godt og Ondt, men uden videre
„slaaer det Ondes Mysterium sammen med alle øvrige, ikke
blot overnaturlige, men ogsaa naturlige Mysterier.“ (Jfr.
Dogm. Oplysn. S. 65); ja --- saa forvirret er denne paatrængende,
fanatiske Prof. N., at han, skjøndt aabenbart
betagen af Iver for en eensidig „Intellectualisme, Kriticisme
og Scholasticisme“, dog ikke desto mindre udgiver sig
for Troens Forsvarer (Dogm. Oplysn. S. 71), mener (Oplysn.
S. 72) at tale i Troens Eenfold om de guddommelige
Tings Ubegribelighed, er begyndt at blive „absolut i Troen“,
og „har begyndt at tale med den forunderlige Parrhesie og
Vished, som er forbausende for andre simple Troende, der
ikke saa pludselig kunne hæve sig til de høieste Troeshøider.“
--- Med Alt dette forudseer nu Dr. Martensen (Oplysn.
S. 76), „at den Tid ikke vil være fjern, da der atter
vil foregaae en Hovedforandring i Prof. Ns Anskuelser; thi
Dr. Martensen kan umulig antage, „at en saa levende, saa
kraftig og, tiltrods for al sin Flygtighed og dilettantmæssige
Ustadighed, dog væsentlig grundig Natur længe kan finde
sig tilfredsstillet ved at gaae den Vei, han nu er slaaet ind.“

See, dette er uden Fordreielse og Overdrivelse den Belysning,
hvori Dr. Ms dogmatiske Oplysninger stille dem,
der have tilladt sig at rokke ved ποῦ στῶ, og betvivle
% --- p. 9 ---
\setcounter{footnote}{0}%
\opage{9}Videnskabeligheden af hans „christelige Dogmatik.“ Men ---
er det martensenske ποῦ στῶ nu ogsaa virkelig saa urokkeligt,
og den Sikkerhed, Overlegenhed og Myndighed, hvormed
han optræder i sine dogmatiske Oplysninger, saa afgjørende
og fast begrundet, at Enhver, der vover at tvivle
derom, uden videre maa finde sig i at blive anseet for et
forvirret, fremfusende, paatrængende, fanatisk Hoved? Dette
Spørgsmaal er dog vel neppe uden al Betydning. Vi faae
nu at see.

Før jeg imidlertid gaaer nærmere ind paa Sagen selv,
maa jeg foreløbig henlede Læserens Opmærksomhed paa to
Punkter, nemlig: det besynderlige Forhold, hvori Forfatteren
af de „dogmatiske Oplysninger“ stiller sig til Pseudonymen,
Johannes Climacus, og det paafaldende Hastværk,
hvormed han affærdiger Mag. Stilling.

Med megen Bestemthed afviser Dr. Martensen ethvert
Hensyn til de af Mag. Kierkegaard udgivne Skrifter. Han
erklærer saaledes (S. 12): „Hvad nu angaaer Johannes
Climacus og de øvrige Pseudonymer, da ere disse i nærværende
Sammenhæng mig aldeles uvedkommende.“ Men det
lader sig dog alligevel ikke negte, at den tilsyneladende Bestemthed,
hvormed den høiærværdige Forfatter viser Pseudonymen
fra sig (en Bestemthed, der strax efter, S. 13, udgives
for Beskedenhed) forraader en vis indre Forlegenhed,
som end ikke den dogmatiske Sufficance, hvormed Dr. M.,
her, som ellers overalt, fremsætter sine Paastande, kan skjule.
Det Uddrag af Pseudonymen, der udgjør den første Halvdeel
af min „undersøgende Anmeldelse“, er udentvivl kommet
Hs. Høiærv. lidt ubeleiligt, just fordi det kom --- fra den
Kant. Dr. M. har nemlig af visse Grunde allerhelst ønsket
% --- p. 10 ---
\setcounter{footnote}{0}%
\opage{10}at undgaae ethvert umiddelbart Sammenstød med en vis
Johannes Climacus og den hele „vidtløftige Litteratur“,
hvortil denne Skikkelse hører; men, da han nu saae, at
den Basis, hvorpaa min kritiske Vurdering af hans dogmatiske
Speculation hviler, var et aabenlyst Uddrag af Pseudonymen,
har han, som det synes, følt, at han ikke vel kunde
indlade sig paa min „undersøgende Anmeldelse“, uden at sige
et Par Ord om Pseudonymen.

Saa kom Dr. Martensen da omsider i et Slags
umiddelbar Berøring med Pseudonymen; men --- er det
ikke mærkværdigt! --- neppe faaer Hs. Høiærv. sin „ligefremme
Meddelelse“ begyndt med et Par Ord, før det ryster
i ποῦ στῶ, og her viser sig en indre Usikkerhed, en Tvivlraadighed
og Vaklevorrenhed, der ellers er den martensenske
„Genius“ saa fremmed.

Det ubehagelige vis à vis med den polemiske Pseudonym
gjør jo virkelig den ellers bestemte Dogmatiker saa
tvivlraadig og vankelmodig, at han end ikke er enig med
sig selv om, hvorvidt han vil vedgaae at have læst Pseudonymen
eller ikke læst ham, forstaaet ham eller ikke forstaaet
ham: han vil, som vi nu skulle see, baade det Ene og det
Andet, og dog alligevel hverken det Ene eller det Andet,
en Tilstand, der kategorisk udtrykker Vankelmodighed.

Paa den ene Side vil Dr. Martensen nyde Fordelen
af sit fine Bekjendtskab med Mag. Kierkegaards Skrifter,
og giver ikke utydelig at forstaae, at han har en langt
anden (og da naturligviis ogsaa langt rigtigere) Opfattelse
af disse Skrifter, end jeg, der dog vedgaaer at have studeret
dem; idet han jo (S. 12) antager, „at disse Skrifter
bevæge sig i en heelt anden Retning, end den af Prof.
% --- p. 11 ---
\setcounter{footnote}{0}%
\opage{11}Nielsen indslaaede.“ --- Paa den anden Side vil Dr. Martensen
slet ikke lægge nogen Vægt paa sin Opfattelse af
Pseudonymerne, „thi“ --- siger han (S. 13) --- „mit Kjendskab
til denne vidtløftige Litteratur er, som sagt, kun saare
ringe og fragmentarisk.“

Da nu Dr. Ms Kjendskab til Pseudonymernes „vidtløftige
Litteratur“, som sagt, er „saa ringe og fragmentarisk“,
bliver det Uddrag af Johannes Climacus, hvoraf Hs.
Høiærv. saa høilig generes, med dogmatisk Bestemthed erklæret
for et Indbegreb af Sætninger, for hvilke Ingen uden
jeg skal være ansvarlig.

Det Ansvar, der i denne Henseende paaligger mig, skal
jeg da heller ikke fragaae: jeg bør virkelig staae til Ansvar
for, at de omhandlede “Sætninger“ ere Pseudonymens egne,
og at de ikke i noget Væsentligt ere forvanskede af mig; derimod
er der noget Andet, hvorfor jeg ikke tør staae til Ansvar:
jeg tør nemlig ikke svare for, hvorvidt Dr. M. vil vedgaae at
have forstaaet disse „Sætninger“, eller ikke. --- Da det jo fornemmelig
er disse „Sætninger“ paa hvis eftertrykkelige Bekjæmpelse
og grundige Gjendrivelse de „dogmatiske Oplysninger“
gaae ud: skulde man rigtignok formode, at Dr. M.
ligefrem vilde vedkjende sig at have forstaaet „Sætningerne.“
Men ganske vis derpaa kan man dog alligevel ikke være; thi S. 17—18 beklager Hs. Høiærv. sig høilig over, at hans
Recensent ikke har forsøgt sig i en selvstændig Fremstilling,
„for“ --- saa lyde hans Ord --- „at opklare disse for mig
og for saamange Andre uforstaaelige Sætninger, men\dots{}.
har affærdiget os med Pseudonymens egne Ord, hvorved
vi bleve ligesaa kloge, som vi vare tilforn. Thi naar det
gik os saaledes med det grønne Træ eller Pseudonymens
% --- p. 12 ---
\setcounter{footnote}{0}%
\opage{12}Skrifter, at vi ikke kunde faae det rette Syn derpaa, hvorledes
skal det da gaae os med det tørre, eller med Prof.
Nielsens Excerpt?“

Altsaa --- hvis jeg ikke feiler --- vedgaaer Dr. M., at
han endnu ikke har forstaaet Pseudonymens „Sætninger.“
Ja han vedgaaer endnu Mere, thi han vedgaaer (S. 18),
at han tilligemed „saamange Andre“ virkelig engang har
taget sig for at studere Pseudonymerne og arbeide sig ordentlig
igjennem den hele „vidtløftige Litteratur“, men har
maattet opgive det, være sig paa Grund af hans „individuelle
Aandsretning“ eller anden Forhindring.

Dr. M., der jo netop lægger saamegen Vægt paa den
grundige Forskning, Dr. M., som jo netop i sine „dogmatiske
Oplysninger“ har leveret os saa værdifulde Prøver paa
sine lærde Forskninger i Tertullian og Anselm, den forskende
Dr. Martensen kan dog umulig beklage sig over ikke at
kunne forstaae en original Forfatter, naar han ikke vil gjøre
sig den Uleilighed at forske, --- „\textit{negligentiæ mihi videtur};
ja, \textit{negligentiæ mihi videtur}!“ Dr. M. har altsaa for Alvor
forsket i Pseudonymerne, forsket efter Mening og „sammenhængende
Erkjendelse“, men --- hans alvorlige Forskning
har været forgiæves.

Dog, hvad siger jeg, Dr. M. har jo slet ikke forsket
i Pseudonymerne, ikke engang gjort sig den Uleilighed at
læse dem ordentlig igjennem; han tilstaaer det selv: „mit
Kjendskab til denne vidtløftige Litteratur er, som sagt, kun
saare ringe og fragmentarisk.“ --- \textit{negligentiæ mihi videtur},
ja, \textit{negligentiæ mihi videtur}! --- men han tilstaaer det jo
selv, og man maa troe en saa alvorlig Mand paa hans Ord.

Dr. M. kan altsaa erklære Pseudonymerne („det grønne
% --- p. 13 ---
\setcounter{footnote}{0}%
\opage{13}Træ“) for uforstaaelige, endskjøndt han egentlig slet ikke har
læst dem; Dr. M. kan altsaa nok gjendrive Pseudonymens
„Sætninger“, men han kan ikke forstaae dem: see det er
Punkter, for hvilke jeg ikke vil være ansvarlig\footnote[1]{Dr. Martensen beklager sig over, at han ikke kan forstaae Pseudonymerne, men kun i disse Skrifter finde, hvad han i Forordet til sin Dogmatik kalder „Strøtanker, Aphorismer, Indfald og Glimt“. Denne Klage afvises som ubeføiet. Naar man kun læser en Forfatter flygtig og „fragmentarisk“, hvorledes skulde man saa kunne finde Mere end Strøtanker og Aphorismer? Den „ligefremme Meddelelse“ af Dr. M., at han ikke har kunnet forstaae Pseudonymerne, bliver altsaa her at opfatte som „en indirect Meddelelse“, at nemlig Hs. Hv. har forsømt at læse.}.

Dog, om Pseudonymens „indirecte Meddelelse“ saavel
som om Dr. Martensens og mit Forhold til samme skal jeg
siden hen udtale mig lidt udførligere; jeg gaaer nu over til
det andet Punkt: Behandligen af Mag. Stilling. % PRINTED AS IS: Behandligen (for Behandlingen)

Blandt de Ankeposter, Dr. M. fremfører imod min
„undersøgende Anmeldelse“, udhæves især den, at jeg ikke
skal have seet mig istand til at gjengive Pseudonymernes
Problem med mine egne Ord, eller gjøre tilbørlig Rede
for, hvorledes den foregivne Splid imellem Tro og Viden
egentlig opkommer; men derimod har paa en ret discipulær
Maade confunderet den indirecte og ligefremme Meddelelse
med hinanden.

Dersom Dr. M. nu virkelig for Alvor lægger nogen
Vægt paa denne mod mig fremsatte Klage, dersom det nu
virkelig er hans oprigtige Ønske at faae Problemets Tilblivelse
beskreven, omskreven og tydeliggjort i en ligefrem
og docerende Form: maa man da ikke høilig undre sig over,
at Hs. Hv. med en saa yderlig, eller rettere, med en saa inderlig
Indignation har kunnet forkaste Mag. Stillings
Bi%
% --- p. 14 ---
\setcounter{footnote}{0}%
\opage{14}drag. Alt hvad den høiærværdige Dogmatiker saa høilig
savner hos mig, vil han jo netop kunne finde hos Mag.
Stilling.

Naar Dr. M. altsaa opholder sig over min store Udygtighed,
at jeg nemlig er „saaledes indviklet i Svøbet af den
nye Viisdom, og endnu bevæger mig saa besværligt, saa
discipulært i denne Viisdoms første Elementer, at jeg ikke
engang formaaer at foredrage ham mine Læresætninger med
mine egne Ord“; naar Dr. M., siger jeg, fra sin gamle
Viisdoms ophøiede ποῦ στῶ seer ned paa min Ringhed og
overskuer min Uformuenhed: da burde Hs. Høiærv. dog idetmindste
forunde Mag. Stilling et mildere, mere indbydende
Blik, og glæde sig over, at Magisteren i saa Henseende er
langt mere fremmelig, end jeg.

Mag. Stilling er nemlig ikke, som jeg, saaledes indviklet
i Svøbet af den nye Viisdom, at han skulde have
tabt Brugen af sin egen Tunge, thi han formaaer virkelig
at foredrage Dr. Martensen „sine Læresætuinger med sine % PRINTED AS IS: Læresætuinger (u for n)
egne Ord“; Mag. Stilling har da heller ikke, som jeg, det
Uheld at confundere den indirecte Meddelelse med den directe,
thi hans Meddelelse er saa ligefrem og docerende,
som muligt. Da jeg imidlertid til min store Forbauselse
seer, hvorledes Dr. M. har lagt Mag. St.s Skrift tilside,
for ikke mere at optage det, vil jeg, for at bevise, hvad jeg
siger, ved denne Leilighed optage dette Skrift igjen, og i faa
Hovedtræk erindre Læseren om dets Indhold.

I sin med logisk Skarphed udarbeidede Afhandling behandler
Mag. St. netop de to „Læresætninger“, hvorpaa det
i denne Undersøgelse især kommer an, nemlig I) \textit{Quæro intelligere},
\textit{ut credam}, og II) \textit{Credo, ut intelligam}. Altsaa:

% --- p. 15 ---
\setcounter{footnote}{0}%
\opage{15}„I) \textit{Quæro intelligere, ut credam} (Jeg søger ved
Tanketilegnelse, ved videnskabelig Erkjendelse, videnskabelig
Beviisførelse at gjenoplive min ved Reflexion lidt afmattede
Tro).“

Under denne Overskrift afhandler Mag. Stilling („Om
den indbildte Forsoning af Tro og Viden“, S. 3—44) det
fleersidige Forhold, hvori den menneskelige Tænkning kan
stille sig til sin Gjenstand, og beviser saa deraf, at det religiøse
Object, Troens og Haabets Gjenstand, deri er forskjelligt
fra ethvert andet Tankeobject, at det, naar Spændingen
naaer sit Høieste, endog frastøder og tilbageskrækker den efter
Viden higende Subjectivitet. Dette frastødende Forholds
dybere Natur er imidlertid saa langt fra at blive den Reflecterende
tydeligt ved første Øiekast, at tvertimod den i
Theorierne arbeidende Menneskeaand har maattet gjennemgaae
mange og smertelige Erfaringer, for den lærte at indsee,
hvad dette vil sige. Thi den af Troen befrugtede Tænkning begynder egentlig med at betragte det christelige Object
som ethvert andet Object, der dog alligevel efterhaanden
maa lade sig meer og meer forarbeide og tilegne i videnskabelig
objectiv Viden.

„Lad det Christelige endog“ --- saaledes lader Forf.,
S. 7—8, Aanden tale i Monolog --- „kalde sig selv „en
haard Føde“, min Tankekraft, der er forsynet med saa stærke
Fordøielsesorganer, skal vel nok fordøie det; ja Opgaven, at
tænke det, skal og maa lykkes; thi først saaledes er det
egentlig mit, først saaledes er det egentlig inderliggjort,
først saaledes er jeg ligeoverfor det en fri uendelig Subjectivitet,
befriet fra al udvortes Autoritets utaalelige Tryk;
først saaledes er mit Forhold til det christelige Object et
% --- p. 16 ---
\setcounter{footnote}{0}%
\opage{16}Forhold, der er den menneskelige Aand værdigt, thi --- „der
Mensch, da er Geist ist, darf und soll sich selbst des Höchsten
würdig achten“, von der Größe und Macht seines Geistes
kann er nicht groß genug denken; und mit diesem Glauben
wird nichs so spröde und hart seyn, das sich ihm nicht eröffnete.
Das zuerst verborgene und verschlossene Wesen des
Universums hat keine Kraft, die dem Muthe des Erkennens
Widerstand leisten könnte; es muß sich vor ihm aufthun,
und seinen Reichthum und seine Tiefen ihm vor Augen legen
und zum Genusse geben.“

Det er Hegels Aand, der taler; det er den theoretiske
Videns Særkjende, at den altid forudsætter et ligefremt
Forhold imellem det erkjendende Subject og Erkjendelsens
Object, et Forhold, hvori Subjectiviteten tilsidst maa blive
det Overgribende. I dette Punkt er der altsaa ikkun den
Forskjel imellem den hegelske og den theologiske Speculation,
at den hegelske Tænkning ikke bliver staaende paa Halvveien
men fuldfører Conseqvensen. Først ved dette dristige Forsøg
blev det da tilfulde aabenbart, at det christelige Object,
Troens og Haabets Gjenstand, er ueensartet med Videnskabens,
og at der er en Uendelighedens Forskjel mellem det, at
hvile i objectiv Vished, og det, at vandre i Tro og i Haab.
Forsøget viste sig (Anf. Skr. S. 12) som „et Overgreb, mod
hvilket Guddommen ligesom protesterer med de Ord: „nei!
her er ikke Du, Mennesket, men jeg, Herren, den „Overgribende“,
der gaaer ud over dit Begreb, et Kors for din
Tanke, et Kors, der kun kan bæres af dit Hoved ved at
tage alle Tanker tilfange under Troens Lydighed.“

Denne „Troens Lydighed“ er imidleitid ikke at forvexle % PRINTED AS IS: imidleitid
med den indsmigrende Underdanighed, hvormed en mystisk
% --- p. 17 ---
\setcounter{footnote}{0}%
\opage{17}uklar Tænker kan faae isinde at lade (S. 14) „sin Tankekraft
engagere som Dogmernes Jabroder; et Engagement,
der vel stundom tilbydes den af dogmatiserende Tænkere.
Ligesom nemlig en Jabroder stedse engageres efter den Methode,
at han tilsyneladende og for et Skins Skyld beholder
en fri selvstændig Mening, medens han dog i Virkeligheden
kun skal snakke Herren efter Munden: saaledes kan man
ogsaa stundom falde paa at ville engagere Tankekraften
paa samme nedværdigende Betingelser, og derved opnaae
dens servile Underdanighed, en Underdanighed og et Forhold,
der rigtignok er heelt forskjelligt fra det Forhold,
hvor Tanken, efter ridderlig Kamp, omsider giver sig til
Fange under Troens Lydighed.“

Under den fortsatte Anstrengelse for Tilegnelsen af
Aabenbaringstankerne kunne nu vistnok „mangfoldige Skuffelser
og Nuanceringer i det reflecterende Subjects Forhold
til Objectet“ vise sig; „men“ --- hedder det S. 15 --- „væsentlig
er der i dette Forhold kun tre Former mulige, alt
eftersom enten I. overveiende Objectets (de saakaldte Aabenbaringstankers)
Selvstændighed hævdes; eller II. overveiende
det tænkende Subjects Selvstændighed hævdes ; eller III. baade
Objectets og Subjectets Selvstændighed hævdes samtidigen.
Under I. henfører Mag. St. det Slags Tænkere , der paatage
sig „at tænke og at fremstille som tænkelig hele den
saakaldte gamle orthodoxe Christendom“ (Göschel).

Standpunktet II. repræsenteres af Hegel og Marheineke.
Paa dette Standpunkt er „Spændingen mellem Object
og Subject egentlig bortfalden.“ Det christelige Object
blev her egentlig nedsat til et upræget modstandsløst Stof,
der maatte finde sig i at corrigeres efter Tankens Fordring.
% --- p. 18 ---
\setcounter{footnote}{0}%
\opage{18}„Hvad der (S. 20) fremfor Alt var fornødent, det var en
Kritik, som gjorde Ende paa Skinforeningen, (som kun var
opnaaet ved Subjectets Overgreb over Objectet), og som
fremstillede det Christelige i sine skarpe Conturer, saa at
det ret maatte vorde Subjectet iøinefaldende, at Objectet
(Dogmet) vilde være en haard Føde for dets Tanke.“

Først paa Standpunktet III. er Partiet lige. „Objectet
er i denne Form af Forholdet meer end Stof, Subjectet
meer end Jabro'er. Og, idet saaledes begge Leds Selvstændighed
samtidigen hævdes, maa det kunne komme til
Afgjørelse, om det selvstændige Subject formaaer ved Erkjendelsesproces
at tilegne sig det uforvanskede Object eller
ikke.“ --- Hvorledes „Spændingen“ her opkommer, og hvad
den væsentlig betyder, ere nu de Cardinalpunkter, Mag.
St. behandler med ligesaamegen Varme, som indtrængende
Skarphed.

Mag. St. opfatter „Spændingen“ fra to Sider: den
theoretiske og den praktiske. „I det theoretisk-spændte Forhold
(S. 23) beroer Spændingens Intensitet og Omfang
paa Intensiteten og Omfanget af Subjectets Reflexionskraft.
Jo intensivere Reflexionskraft, desto stærkere theoretisk Spænding;
og det er derfor ingenlunde ethvert Subjects Sag at
komme i denne Spænding. I det personligt spændte Forhold
derimod er væsentlig ethvert, selv det simpleste Menneske,
saasnart han af Tilværelsen saaledes udkaares til
Trængsel og Smerte, at hvad der --- seet og opfattet paa
menneskelig Viis --- af ham, den Paagjældende, maa føles
og opfattes som den dybeste Ulykke og Modgang, der kuer
og hæmmer hans Liv, dog, seet og opfattet paa christelig
Viis, skal være det Bedste, Guds Styrelse, Medgang. Men
% --- p. 19 ---
\setcounter{footnote}{0}%
\opage{19}først Foreningen af det theoretisk spændte og det praktisk
(personlig) spændte Forhold i eet Subject danner Spændingens
Høieste og den Livs-Situation, i hvilken det reflecterende
Subject maa kunne faae lidt at vide om, hvad det
er, der egentlig bydes ham, naar det bydes ham at troe.“ ---
Den Læser, der med sund Alvor og stille Eftertanke vil
prøve, hvad Afhandlingen (S. 24--- 44) meddeler om den
christelige Forsynstro og det christelige Bønsforhold, skal i
Sandhed „faae Syn“ paa Problemets Tilblivelse og opdage
et „Paradox“, der ingenlunde lader sig tillæmpe til dogmatisk
Speculation derved, at man søger „det tilsvarende Udtryk
i Modersmaalet.“ --- „II) \textit{Credo, ut intelligam} (Jeg, en Reflecterende, er
fra Reflexionsveiene heldigviis sluppen tilbage til et fast,
urokket bibelsk-kirkeligt \textit{Credo}, og søger nu atter fra samme
af meer og meer at hæve min Troes-Vished til Intelligensens
Vished).“

Under denne Overskrift afhandler altsaa Mag. Stilling
(S. 45—92) den Sætning, som vel med Rette maa kaldes
Dr. Martensens Cardinalsætning, den Sætning, som egentlig
betegner vor speculative Dogmatikers ποῦ στῶ; hvoraf
det da foreløbig kan sees, at Dr. M. gjør sine Recensenter
stor Uret, naar han beskylder dem for, at de slet ikke ville
sætte sig ind i hans „Grundanskuelse“, og begynder sine
„dogmatiske Oplysninger“ med at udslynge Luthers Ord:
„Ihr habt einen andern Geist als wir!“ --- imod dem.

Den udførlige Drøftelse af „\textit{Credo, ut intelligam}“ har
Mag. St. allerede indledet i sit foregaaende Afsnit; idet
han (S. 37) gjør opmærksom paa, at det nemlig slet ikke
gaaer „bedre med Reflexionen \textit{post fidem} (ɔ: efterat jeg har
% --- p. 20 ---
\setcounter{footnote}{0}%
\opage{20}taget Troespartiet) end \textit{ante fidem} (ɔ: førend jeg tog Troespartiet,
men endnu stak lige midt i Reflexionsmodsætningerne).
Reflexionsmodsætningen (Vantro) blev ved første
Troesskridt overvunden, ikke ved Reflexionen, men \emph{ved} og
i troende Leven, der dicterede Lydighed.“ Det Samme
vil og maa nu skee ved ethvert af de følgende Skridt.
Derfor er da ogsaa (S. 38) „visse dogmatiserende Tankeres
Distinction mellem et speculativt Beviis \textit{ante fidem}
(der skal være umuligt) og et speculativt Beviis \textit{post fidem}
(der skal være muligt) en fuldkommen falsk og uholdbar
Distinction.“ --- Troen er, med andre Ord, Qvalitet, og
dens personlige Overbeviisning qvalitativt forskjellig fra Reflexionen,
der altid har ikke blot et Vidnesbyrd \textit{pro}, men
ogsaa et Vidnesbyrd \textit{contra}.

Dersom nu Dr. Martensens Dogmatik virkelig var forsvarlig
bygget og videnskabelig grundet paa Sætningen:
\textit{Credo, ut intelligam}, da maatte den, ifølge Mag. Stillings
Paastand, saavel i sit almindelige Anlæg, som i Behandlingen
af de enkelte Dogmer, afgive et klart og fuldstændigt
Vidnesbyrd om, hvorledes Troen (credo) og Reflexionen
(\textit{ut intelligam}) egentlig forholde sig til hinanden. Dogmatikeren
burde altsaa i Gjerningen vise, at han baade kjender
Troens Qvalitet samt det \textit{genus cognoscendi}, Troen
kræver, og tillige kjender den objective Videns (Intelligensens)
Qvalitet samt det Resultat, hvortil en fri og alsidig
Reflexion nødvendigviis maa føre. Men just i denne Henseende
mener Magister Stilling, at det ikke hænger rigtig
sammen med Dr. Martensens „sammenhængende Erkjendelse.“

Magister Stilling anker paa, at man ikke kan see,
hvad det er, Dr. Martensen paa sin Reflexionsvandring fra
% --- p. 21 ---
\setcounter{footnote}{0}%
\opage{21}Troen til Intelligensen og fra Intelligensen til Troen egentlig
har erfaret.

„Erfarede han“ --- spørger Recensenten (S. 50) --- „at
han væsentligen og hovedsageligen ikke kunde tænke det Christelige
?“

Dette Spørgsmaal besvares benegtende. „Vel“ --- hedder
det S. 53 --- „siger Prof, M. (Dogm. S. 108), at hans
forsøgte Viden, hans speculative Erkjendelse, ikke er ganske
adæqvat, at der i den høieste speculative Viden er et Moment
af Ikke-Viden. Men for det Første er dette Udtryk
kun et afmattet og høist ucorrect Udtryk for Tankekorset,
for Anstødet, Paradoxet og Reflexionsvidnesbyrdet \textit{contra};
og dernæst, hvis saa Udtrykket „uadæqvat Viden“ ved nærmere
Forklaring ikke kommer til at betyde Andet, end en
Viden, der, qua Viden, egentlig slet ikke fortjener Videns
Navn: saa bliver jo egentlig hele Resultatet af Vandringen
paa den „menneskelige Erkjendelses Vei“, hele Resultatet af
den forsøgte Viden en forsøgt Ikke-Viden, der tilbeder Guds
forborgne Viisdom. Men heelt og holdent og uforbeholdent
at vedgaae en Ikke-Viden, eller bedre (thi det er den
correcteste Formel), at man, ved at reflectere tilbunds, reflecterer
sig ind i de dybeste og forfærdeligste Reflexionsmodsætninger
af \textit{pro} og \textit{contra}, og støder paa Tankeforset: det % PRINTED AS IS: Tankeforset (f for k; three readers)
vægrer og krymper vor Dogmatiker sig for.“

„Erfarede han da“ --- spørger Recensenten videre (S.
62) --- „at han væsentligen og hovedsageligen kunde tænke
det Christelige?“

Ogsaa dette Spørgsmaal maa besvares benegtende.
„Det vil ved denne Prøvelse“ --- siger Recensenten (S. 66 ---67) „ret øiensynligen vise sig, at vor Dogmatikers
% --- p. 22 ---
\setcounter{footnote}{0}%
\opage{22}„\textit{Credo}“ i Gjerningen er et scholastisk hemmende \textit{Credo},
der lægger Reflexionsfriheden under Censur, og ikke tilsteder
den andre Yttringer, end saadanne, der ere Dogmatikerens
Øre velbehagelige; en høist besynderlig Maade at betjene
sig af, for at adspørge og udfritte Tanke-Autoritetens Mening
om Dogmet.“

For nu nærmere at eftervise, hvorledes Dr. Martensen
har afbenyttet sit credo, \textit{ut intelligam}, udhæver Mag.
Stilling forskjellige Dogmer, og kommer saa hver Gang
til det Resultat, at Dogmatikeren, der ikke vil blive staaende
ved \textit{Ita}, men skride frem til \textit{Quare}“, ikke paa noget Punkt
holder, hvad han lover, men ved hvert Dogme lader, som
han tænker, og tænker ikke endda.

Dette sees nu for det Første i Behandlingen af Trinitetsdogmet
(S. 68---81). Allerede i Begyndelsen begaaer Dr.
M. den Inconseqvens, at ville bevise Trefoldigheden i Guds
ontologiske Selvaabenbarelse, uden dog at ville bevise Guds
Tilværelse. „At bevise“ --- hedder det, Dogmat. S. 89 ---
„den aabenbarede Guds Tilværelse kan ikke være Opgaven.“
Hvortil Mag. St. bemærker: „Kan det være Opgaven for
Dogmatikeren at bevise, at den aabenbarede Guds Personlighed
maa tænkes som en trefoldig Personlighed; hvorfor kan
det saa ikke være Opgave at bevise den aabenbarede Guds
Tilværelse, eller bevise hans Personlighed. Thi at bevise hans
Tilværelse er just at bevise hans Personlighed.“ --- Men heller
ikke den Prøve, Dr. M. her aflægger paa sin Begribeliggjørelse
af Guds Treenighed, indeholder mindste Bidrag til at
føre denne Opgave videre ad Speculationens Vei. Vistnok
siger Dr. Martensen (Dogm. S. 128): „At tænke Væsenstriniteten
ontologisk er at tænke den nødvendige Grundform
% --- p. 23 ---
\setcounter{footnote}{0}%
\opage{23}for Guds personlige Liv, at tænke de Momenter i Guds
Væsen, uden hvilke Personlighed og Selvbevidsthed er utænkelig.
Men Mag. Stillings Kritik gjør det kun altfor indlysende,
at den høiærværdige Dogmatiker aldeles ikke har
tænkt Noget af det, der her egentlig skulde tænkes. De velbekjendte
„Momenter i Guds Væsen“, der skulde afgive Treenighedsbegrebets
Grundlag, ere laante fra Tænkere, hvis
Navne Dr. M. end ikke har nævnet. Dette fra Speculationen
laante Trinitetsbegreb viser sig her endnu ligesaa
ueensartet med den christelige Treenighedslære, som i tidligere
speculative Dogmatiker; hvorfor ogsaa det gamle
Problem, hvorvidt der maa tænkes at være eet eller tre
Jeg'er i Guddommen, aldeles ikke løses i den Martensenske
„Speculation“, der blot hjælper sig ved at neddysse Reflexionen
og fortie Vanskelighederne.

Paa lignende Maade vedbliver Mag. Stilling sin Kritik
over den speculative Grundighed, hvormed Dr. Martensen
behandler de øvrige Dogmer, nemlig Incarnationsdogmet
(Anf. Skr. S. 81 ---85), Arvesyndsdogmet og den personlige
Tilregnelighed (S. 81 --- 88), Udødelighedsdogmet
(S. 88 --- 92). At Kritiken over disse enkelte Dogmer er
noget kort, maa man vistnok forklare deraf, at den ærede
Recensent ikke har havt „umiddelbar og levende Ideeanskuelse“
nok til ret at dvæle ved det positive Apparat af
Skriftsteder eller ved den geniale Forfatters aandrige Omskrivninger
af det Overleverede, men derimod har uden Omstændigheder
blottet de egentlige Knudepunkter, for derved
at godtgjøre, at Dr. M. i Grunden slet Intet har tænkt,
af hvad her dog væsentlig burde tænkes.

Hvorvidt det for Øvrigt er noget videre bevendt med
% --- p. 24 ---
\setcounter{footnote}{0}%
\opage{24}den her skizzerede Kritik vil jeg imidlertid helst overlade til
Læseren selv at betænke; men derimod kan jeg ikke undlade
at udtale min Forundring over Dr. Martensens paafaldende
Aversion for det Stillingske Indlæg.

Dr. Martensen vil nu, som sagt, Intet vide af Mag.
Stillings Kritik; men, efter min ringe Mening, burde Hs.
Høiærv. fra sit Synspunkt just have lagt betydelig Vægt
paa den Stillingske Kritik, og, naar han fandt, at han i
ethvert Tilfælde kun burde indlade sig med Een, da hellere
have valgt Mag. Stilling end mig; Hs. Høiærv. burde,
hvis jeg ikke feiler, aabenbart have gjort det af to Hensyn:
deels for sin egen Skyld, deels for Folks Skyld.

Hvad det første Hensyn angaaer, klager jo Dr. M. over,
at han ikke ret kan forstaae „den indirecte Meddelelse“; saameget
mere burde han da have paaskjønnet, at Mag. Sts
Skrift ved sin „ligefremme Meddelelse“ just vilde være ham
en god Veiledning til at komme ordentlig ind i Sagen, og
see, hvad det var, man her udgav for „Problemet.“ Thi
Problemet ligger, som Mag. St. viser, ganske rigtig deri,
at Troesobjectet staaer i et ganske andet Forhold. til den
menneskelige Erkjendelse, end et hvilketsomhelst af de Objecter,
der ellers blive Gjenstand for objectiv Viden.

Havde Dr. M. nu fæstet sin Opmærksomhed paa dette
Troesobjectets \emph{omvendte} Forhold til det vidende Subject,
havde Dr. M. indladt sig paa denne Vanskelighed, istedetfor
at hjælpe sig med pure Forsikkringer, og „fremdeles“ vedblive
(Dogm. Oplysn. S. 55) „at betragte Dogmatikerens
Forhold til Aabenbaringen paa tilsvarende Maade, som Naturforskerens
til Naturen“; havde Dr. Martensen Punkt for
Punkt viist det Intetsigende i Mag. Stillings Argumenter
% --- p. 25 ---
\setcounter{footnote}{0}%
\opage{25}(især det med „Jabroderen“ og „Rækken \textit{contra}“); ja, havde
Dr. M. virkelig følt sig oplagt til at parere det Stillingske
Angreb med den Sikkerhed, Rolighed og Overlegenhed, som
Bevidstheden om en retfærdig Sag giver: her havde været
vundet Meget. Ikke blot var Dr. M., som sagt, ved denne
Fremgangsmaade kommen (hvad han nu desværre ikke er)
ordentlig ind i Sagen; men det kunde da maaskee ogsaa være
lykkedes ham endnu bedre at „orientere“ sine Læsere „over
Sagens virkelige Stilling.“

Ja havde Dr. M. blot kunnet nedlade sig til at gjendrive
Mag. Stillings kritiske Indsigelser: her havde visselig
været vundet Meget.

Alle velvillige Læsere (især alle de Læsere, der ere saa
velvillige at nøies med Dr. Ms. „dogmatiske Oplysninger“
uden at bekymre sig stort om Modstandernes kritiske „Skriverier“)
vilde da uden Tvivl have raisonneret saaledes: „See
her seer man, hvad det nye Paradoxmagerie har at betyde,
Pseudonymen er jo kun en Humorist, en „Gymnastiker.“ Med
den humoristiske Gymnastiker kan en saa alvorlig Dogmatiker
som Dr. M. ikke indlade sig, da han jo naturligviis ikke
kan vide, om Humoristens „indirecte Meddelelse“ skal være
Spøg eller Alvor. Prof. N. er nu rigtignok ingen Humorist,
men dog en ligesaa letsindig som uheldig Efterligner, der
har forgabet sig i Humoristen og er bleven fremfusende.
Med den „fremfusende“, „paatrængende“ og „letsindige“
Prof. N. kan en saa besindig og grundig Forsker, som Dr. M.
egentlig heller ikke indlade sig; saameget mindre, som Prof.
N. kun giver os „Excerpt“ af Pseudonymen, og ikke er istand
til at udvikle Problemet i „ligefrem Meddelelse.“ Men see,
neppe viser der sig en Modstander, paa hvis docerende For%
% --- p. 26 ---
\setcounter{footnote}{0}%
\opage{26}søg man dog nogenledes kan reflectere, neppe mælder sig
en Kritiker, der dog i nogen Grad formaaer at bruge sin
egen Tunge, neppe optræder Mag. Stilling med sin ligefremme
Meddelelse og sine logiske Forklaringsgrunde, før det
ogsaa viser sig, at vi i Dr Martensen have den Mand, der
ligeledes i ligefrem Meddelelse og ved endnu vægtigere Forklaringsgrunde
er istand til at „feie for Dogmatikens Dør“
og „orientere os Læsere over Sagens virkelige Stilling.“

Naar nu Dr. M. desuagtet lader denne øiensynlige
Fordeel slippe sig af Hænderne og, skjøndt han saa stærkt
opfordrer til „direct Meddelelse,“ dog alligevel afslaaer enhver
Forhandling med den Modstander, der just tilbyder
ham, hvad han ønsker: saa bliver man jo ikke saa ganske
„orienteret over Sagens virkelige Stilling.“

Den selvtænkende Læser (den Læser, der ikke uden
videre tør betragte Dr. Ms. Oplysninger som Troesartikler)
maa da uvilkaarlig spørge sig selv: „Hvad mon der
dog kan have bevæget den høiærværdige Dogmatiker til at
behandle Mag. St. paa en saa haanlig Maade;“ thi, sandt
at sige, her taler jo Dr. M. ingenlunde som den Guds Mand,
der kun tænker for „den hellige almindelige Kirke,“ men
snarere som den dybt krænkede Forfatter, der af en kun altfor
naturlig Selvfølelse forledes til en Ubesindighed: \textit{tantæne}
\textit{animis coelestibus iræ}! Hvad har Mag. St. da syndet, at han
saa brat forstødes af „en Menighedens“ Ven, af Ordets høiærværdige Tjener? I personlige Grovheder og moralske Invectiver
kan Overtrædelsen dog umulig bestaae; thi den paagjældende
Kritik indeholder ingen Yttring, som Enhver i en saadan Sag jo
honnetement kan tillade sig, naar kun hans Argumenter for
Resten holde Stik. Men hvori bestaaer da Mag. Sts. Mis%
% --- p. 27 ---
\setcounter{footnote}{0}%
\opage{27}gjerning? Den skulde dog. vel aldrig bestaae deri, at han
har talt saa frit om Reflexionen som „Jabroder“ og „Rækken
\textit{contra},“ at Dr. M., der slet ikke behøver „Rækken \textit{contra},“
efterdi han alene speculerer udi Gjenfødelsen, deraf har
sluttet, at „denne uværdige Recensent endnu ikke var gjenfødt?“
Misgjerningen skulde dog vel aldrig bestaae deri, at Mag.
St.\footnote[1]{Et Par Spørgsmaal til Prof. Scharling.} har været saa ubeskeden, offentlig at omtale et Factum,
der er af en saa delicat Natur, at Dr. Ms. „Beskedenhed“\footnote[2]{Dogm. Oplysning S. 69.}
ikkun tillader ham at ansee det for „mindre væsentligt,“
det Factum nemlig, at den høiærværdige Dogmatiker
havde, ifølge hans egen udtrykkelige, trykte og offentliggjorte
Vidnesbyrd, et ganske andet Videnskabsbegreb
i Aaret 1839, end det, der nu fra Aaret 1849 ligger
til Grund for hans Dogmatik. Skulde dette være Tilfældet,
saa er Hs. H. vist igjen imod sin Villie kommen til
at give os --- „en indirect Meddelelse.“

For strax at „orientere“ os over „Sagens virkelige
Stilling“, erklærer Dr. M., som vi have seet, at Kritiken
over hans Dogmatik egentlig endnu ikke kan siges at være
begyndt. Jeg kan nu meget godt forstaae, hvad det er for
en Lyseffect, den høiærværdige Forfatter tilsigter ved denne
„dogmatiske Oplysning“, men Betydningen af de Grunde,
hvorpaa han støtter sin Paastand, kan jeg derimod ikke rigtig
forstaae. Han klager over, at Kritiken nærmest har holdt
sig til Indledningen. Men selv om dette ogsaa var Tilfældet,
vilde denne Omstændighed her dog ikke begrunde nogen
Klage. Ja, hvis Kritiken gik ud paa at vise, at Værkets Form
% --- p. 28 ---
\setcounter{footnote}{0}%
\opage{28}og Udarbeidelse ikke svarede til de i Indledningen opstillede
Grundsætninger, saa var det en anden Sag. Men her forholder
Sagen sig omvendt. Kritiken indrømmer tvertimod,
at Udarbeidelsen er saa god, som man ifølge de angivne
„Grundsætninger“ med Billighed kan forlange den\footnote[1]{Her er jo vistnok i denne christelige Dogmatik, seet fra Formsiden, ogsaa noget Behageligt, noget Tiltalende, Noget der i visse Stemninger virker, om ikke just strengt overbevisende, saa dog mildt beroligende. Her er i denne christelige Dogmatik noget Fleersidigt, noget Forsonende, noget Fredsommeligt, noget Føieligt. Her er i denne christelige Dogmatik, naar den betragtes som systematisk Heelhed, tillige noget vist Kunstnerisk, noget Architectonisk, Noget, der paa een Gang ligesom formæler Universitetsbygningen med Slotscapellet, Videnskaben med Religionen: skarp Forstand og naiv Eenfoldighed, moderne Reflexion og mystisk Umiddelbarhed, Begrebets Klarhed og Mysteriets Dunkelhed, Alt smelter sammen i Stilens Eenhed, i festlig Ro, høitidelig Dæmring, som Dagens Skjær med Skinnet af et Alterlys.“ Undersøg. Anm. S. 43.}; kun
Grundsætningernes videnskabelige Holdbarhed er det, Kritiken
modsiger. Da nu Dr. M., hvad han i de „dogmatiske Oplysninger“
ogsaa (S. 5.) selv tilstaaer, har i Indledningen
fremstillet sine Grundsætninger „saa skarpt og bestemt, som
muligt,“ saa indseer jeg virkelig ikke, hvad der kan indvendes
imod, at en Kritik over „Grundsætningerne“ fornemmelig
holder sig til „Indledningen“, og kun deelviis henter
Exempler fra Udførelsen.

Endvidere hedder det (S. 6), at Kritiken nu atter skal
være kommen ind paa „den gamle, brede Vei, der i slet
Bemærkelse kan betegnes som den tydske, hvor man nemlig
tumler sig i lutter propædeutiske Undersøgelser, opvækker
Støvskyer, men aldeles ikke kommer til Sagen.“ --- Men
skulde det nu ogsaa virkelig være aldeles overflødigt, at gaae
ind paa „propædeutiske Undersøgelser,“ og vende tilbage til
% --- p. 29 ---
\setcounter{footnote}{0}%
\opage{29}Principerne, naar man vil udfinde, hvorvidt et dogmatisk
Arbeide har noget Princip eller er --- principløst.

Den Gang Dr. M. skrev sit lille Skrift: „Om Selvbevidsthedens
Autonomie“, og „paa nogle faa Blade“ kritiserede
baade Kant og Schleiermacher og Hegel, var han uden
Tvivl selv af den Mening, at man endda nok kunde prøve
et dogmatisk System uden just at fortabe sig i Detaillen af
det dogmatiske Stof; den Gang ansaae han det endnu ikke
for overflødigt at dvæle ved Indlednings-Spørgsmaalene
og „Grundsætningerne“; men nu, da hans eget dogmatiske
Værk er fuldendt, nu er det anderledes; nu skal der slet
ikke komme Noget ud af at drøfte Dogmatikens Indledning,
nu bliver man strax en „Erasmus Montanus“, naar man
har den „Paatrængenhed“ at ville discutere de dogmatiske
Principer. En betydningsfuld Forandring!\footnote[1]{Dr. Martensen, der netop gjorde sin Lykke ved Universitetet ved at foredrage Standpuncternes Philosophie i kort, populært æstetisk Overblik; Dr. M., der jo netop kom op ved sine „Indledninger om Tro og Viden“ og sine „foreløbige Anstalter til et Systems Opbyggelse“; Dr. M., der i Begyndelsen anmeldte sig som en Ideens Mand, der kun vilde røre Systemernes Nerve: den samme Dr. M. er nu (Dogm. Oplys. S. 5) med Eet bleven saa mæt og træt af alle „theologiske og philosophiske Indledningsspørgsmaal“, at han endog har begyndt at forvandle sig til en Stoffets Mand, der taler om lærde Forskninger i Tertullian. Blot det nu ikke skulde hændes Dr. M., at han over sin dybe Forskning i Tertullian forglemmer Apostlenes Gjerninger, thi at forvexle Festus og Agrippa (jfr. Dogm. Oplys. S. 30 med Ap. G. Cap. 26, 24.), er --- \textit{NB.} for en lærd Forsker --- lidt ubeldigt.} % PRINTED AS IS: ubeldigt (for uheldigt), in the note

At ville, som Dr. M. (Dogmatisk. Oplysn. S. 7.)
udtrykker sig, „affærdige Treenighedslæren, Incarnationslæren,
Læren om Syndefaldet og Arvesynden, om Frihed
og Naade, foruden en heel Cyclus af andre Dogmer paa
% --- p. 30 ---
\setcounter{footnote}{0}%
\opage{30}et Par Blade,“ er vistnok urimeligt, naar der forlanges en
speculativ-dogmatisk Fremstilling af hvert enkelt Dogma for
sig. Men Saameget vil dog vel Dr. M. endnu, naar han
ret besinder sig, sagtens indrømme os, at der er en stor
Forskjel mellem det, at levere en positiv Bearbeidelse af et
Dogma, og det, at underkaste en saadan Bearbeidelse en
videnskabelig Kritik. Det Første vil i Reglen udkræve mange
Blade, det Sidste kan, naar Kritikeren kun rammer det rette
Punct, meget vel skee paa et Par Blade.“ Ogsaa i denne
Henseende havde det været ønskeligt, om Dr. M. ikke blot
havde bladet og talt Bladene i Mag. Sts. Afhandling, men
ogsaa læst og gjennemtænkt, hvad der stod paa disse Blade,
før han saa „strax lagde Hr. Magisterens Skriveri hen, for
ikke mere at optage det.“

Som Beviis for, at der er Noget i, hvad jeg siger,
skal jeg endnu blot tillade mig at citere et eneste Sted af
den Stillingske Bog. Med Hensyn paa Ubiqvitetsdogmet,
der omhandles S. 83, hedder det i en Anmærkning S. 73:
„At Prof. Martensen i Virkeligheden end ikke i fjerneste
Grad formaaer at tænke en punctuel Personligheds Allestedsnærværelse,
røber han da ogsaa ganske naivt i Ubiqvitetslæren,
hvor Talen er om den himmelfarne Gudsson, der
har sit Sæde ved Faderens høire Haand, som er overalt.
„Ville vi“ --- saaledes siger vor Dogmatiker, S. 387 ---
„fastholde og gjennemføre Forestillingen om en ubegrændset
Christusnærværelse i Himlen og paa Jorden, da kunne vi
ikke undgaae at forflygtige Christi Individualitet“ (ja
det indrømmes ligefremt, at en ubegrændset Christusnærværelse
fører til en pantheistisk Opfattelse); „thi selv en forklaret
Individualitet, selv et aandeligt Legeme lader sig ikke
% --- p. 31 ---
\setcounter{footnote}{0}%
\opage{31}tænke uden Begrændsning.“ Her spørger nu Recensenten: „Lader
det sig da i mindste Maade bedre \emph{tænke}, at Christus
er personlig tilstede i hver Nadvere, naar der i een By
eller i 10 a 100 forskjellige Byer samtidigen holdes Nadvere
paa a 100 forskjellige Steder i hver By?“ ---

Denne lille Anmærkning, der ikke optager et Blad, end
sige „et Par Blade,“ er af den Beskaffenhed, at den ikke
blot belyser den Martensenske Behandling af Ubiqvitetsdogmet,
men tillige det principielle Forhold, hvori Dr. M. som
speculativ Tænker har sat sig til de kirkelige Dogmer i
det Hele.

En videnskabelig Recension af vor speculativ-christelige
Dogmatik kan i Henhold til denne Anmærkning kortelig redigeres
saaledes: „En original speculativ Dogmatiker er
opstaaet iblandt os. At denne Dogmatiker baade er original
og speculativ, sees af hans Behandling af alle Dogmer
uden Undtagelse; men at han virkelig er ligesaa original,
som speculativ, og ligesaa speculativ, som original, fremgaaer
i Synderlighed af hans Behandling af Ubiqvitetslæren.
Dette dogmatiske Problem have hverken de ældre eller de
yngre Theologer hidtil kunnet løse paa nogen ret tilfredsstillende
Maade. Først i det foreliggende Værk er det omsider
lykkedes „den dogmatiske Genius“ at løse Problemet og
ophæve „Eensidigheden“ ved en ganske ny Opdagelse. Opdagelsen
er denne: For at være nærværende i Nadveren behøver
den forklarede Christi Legeme endnu ikke egentlig at
være \emph{allestedsnærværende}, det behøver blot at være
\emph{mangestedsnærværende} paa een Gang. Da nu en
saadan Mangestedsnærværelse egentlig maa \emph{tænkes} som en
blot relativ og begrændset Allestedsnærværelse: er her jo
% --- p. 32 ---
\setcounter{footnote}{0}%
\opage{32}Intet iveien for, at en saadan begrændset Allestedsnærværelse
meget vel kan forenes med den forklarede Individualitets
nødvendige Begrændsning. See, det er Problemets Løsning,
en Løsning, hvori der jo ikke er noget Paradox! --- Hvorledes
nu vi Andre, der ikke have kunnet gjøre Opdagelsen, imidlertid
skulle bære os ad med at faae den nylig opdagede
„Mangestedsnærværelse“ virkelig tænkt, siger vor Dogmatiker
ikke: han siger kun, at den kan tænkes og bør tænkes og
maa tænkes; men Mere siger han ikke, og dette er just det
Originale. Hvorledes „Geniussen selv“ bærer sig ad, naar
han logisk skal analysere den „relative Allestedsnærværelses“
Begreb, siger han os heller ikke; han siger bare credo, \textit{ut}
\textit{intelligam}; men, saasnart han bare siger: credo, \textit{ut intelligam},
er det ogsaa, ligesom om det virkelig var tænkt, og
dette er just det Speculative.“

Mon det endnu, for at charakterisere det egentlige
Standpunct, skulde være nødvendigt at gaae ind paa „den
hele Fylde af Dogmer,“ eller kunne vi maaskee lade det
blive ved denne magre Recension, og sige: sat sapienti?

„Dersom den dogmatiske Genius“ --- siger Dr. M.
(Dogm. Opl. S. 54) --- „dengang har været productiv, tør
den vel være productiv endnu,“ og heri maa jeg give ham
Ret. Men det kommer da nu heller ikke alene an paa, at
„en dogmatisk Genius“ er „productiv“; det kommer jo ogsaa
væsentlig an paa, hvad det er, „den dogmatiske Genius“
producerer.

Den dogmatiske Productivitet kan for nærværende Tid
dog vel ikke nærmest gaae ud paa at opdage nye Troesartikler
eller frembringe nye Dogmer. Jeg for mit Vedkommende
kan slet ikke finde, at Dr. Ms. „dogmatiske Genius“
% --- p. 33 ---
\setcounter{footnote}{0}%
\opage{33}har produceret et eneste nyt  Dogme (med mindre det da
skulde være Dogmet om den forklarede Christi „Mangestedsnærværelse“).
Det Indbegreb af Troessætninger, han behandler
i sin Dogmatik, er netop den samme Cyclus af
Artikler, Skriftsteder og Qvæstioner, som længst ere bearbeidede
paa de allerforskjelligste Maader i den lutherske Kirke.
Til Tilegnelsen af dette Stof udfordres hverken noget overvættes
Forraad af „positive Kundskaber“ eller noget særeget
„paa religiøs Livserfaring grundet Bekjendtskab med Aabenbaringens
Sandheder“; man finder det hele Apparat i næsten
ethvert dogmatisk Compendium, ja i nogle Haandbøger (t.
Ex. i Bretschneiders Dogmatik og Hases Hutterus redivivus)
endog langt udførligere, end i Dr. Ms Bog. Den
Martensenske Productivitet har da i dette Stykke ingen lærd
Betydning, som om hans Dogmatik var et nyt og eiendommeligt
Resultat af dogmehistorisk Forskning Betydningen
af Dr. Ms „Productivitet“ maa altsaa her søges i
noget ganske Andet, nemlig i den videnskabelige Dygtighed,
hvormed han heelt igjennem har vidst at opklare Forholdet
imellem Troen og Erkjendelsen.

I denne Henseende fordrer saa Dr. M., at man, for
ret at vurdere hans „dogmatiske Productivitet“, skal (Dogm.
Oplysn. S. 6) „lade Bedømmelsen af Erkjendelseslæren gaae
Haand i Haand med en Prøvelse af selve Dogmeudviklingen,
efterdi disse Tvende gjensidigt oplyse hinanden, og
ikke rettelig kunne forstaaes uden i Forening.“ --- Jeg skal
nu, støttende mig til Mag. Stillings Anmærkning, strax
have den Ære at efterkomme denne Fordring.

„Læren om Christi Sæde ved Faderens Høire --- hedder
det, Dogm. § 177 --- „er udviklet i Læren om Christus,
% --- p. 34 ---
\setcounter{footnote}{0}%
\opage{34}som den, der i Alt gjennemtrængende Nærværelse opfylder
Alt i Alt (τοῦ τὰ πάντα ἐν πᾶσι πληρουμένου).“ --- Her
er altsaa et enkelt, bestemt, positivt Dogme. Til dette Dogme
kan nu Erkjendelsen sætte sig i et dobbelt Forhold, et videnskabeligt
og et religiøst, et objectivt og et subjectivt Forhold.
I begge disse Forhold kan man begynde med Sætningen
\textit{Credo, ut intelligam}; men, som vi strax skulle see, faaer
Sætningen en ganske anden Betydning i den religiøse Erkjendelse
end i den videnskabelige.

I objectivt videnskabelig Forstand gaaer \textit{Credo} ud paa
en historisk Antagelse af et givet Faktum: \textit{Credo}, jeg troer,
ɔ: jeg antager for afgjort vist, at Dogmet i sin bestemte
Form udtrykker Christendommens og Kirkens sande og rette
Mening. \textit{Credo}, jeg antager Dogmet, som det er, \textit{ut intelligam},
ɔ: for at jeg derved kan komme til Erkjendelse.
Men til hvilken Erkjendelse? Til en videnskabelig Erkjendelse,
ɔ: til Erkjendelse af, hvorvidt jeg ad Videnskabens
Vei kan komme til at indsee Dogmets objective Sandhed,
eller ikke. For at afgjøre dette, maa jeg strax henvende
min Opmærksomhed paa to Hovedpunkter, idet jeg nemlig
maa undersøge: a) om Dogmets Indhold lader sig opfatte
i nogensomhelst objectiv Begrebsform, og b) hvorvidt den
dogmatiske Forestillings Realitet lader sig prøve efter nogen
af de os bekjendte objective Love.

I begge Henseender er det imidlertid let at bevise, at
den videnskabelige Erkjendelses Forhold til Dogmet umulig
kan være et positivt, men maa være et negativt Forhold;
thi a) Begreberne: „Allestedsnærværelse“ og „Individualitetsbegrændsning“
ere hinanden modsigende: en allestedsnærværende
Individualitet er utænkelig, og Alt, hvad der
% --- p. 35 ---
\setcounter{footnote}{0}%
\opage{35}kan siges om en individuel Mangestedsnærværelse, --- Nonsens.
Selv om Dogmets Indhold ikke var i saa aabenbar Strid
med Begrebet, som det virkelig er: det vilde dog alligevel
ikke kunne tilegnes i videnskabelig Erkjendelse; thi b) Forestillingen
om en forklaret Individualitet, der „opfylder Alt
i Alt“, hører ikke hjemme i den Tingenes Orden, hvis Love
vi kjende, men søger sin Gjenstand i det Overnaturlige.
Men, saasnart Videnskaben indlader sig med det Overnaturlige,
opløser den sig selv, thi her er ingen Grændse for
Phantasteriet: enhver dybsindig „Genius“, der beraaber sig
paa den „religiøse Livserfaring“, enhver speculativ Charlatan,
der paakalder „den umiddelbare og levende Ideeanskuelse“,
kan da uden videre paastaae, hvad det skal være,
og udgive sine Paastande for videnskabelige Erkjendelser.

Anderledes med den religiøse Erkjendelse. I religiøs
Forstand faaer Ordet \textit{Credo} i Sætningen: \textit{Credo, ut intelligam},
en afgjort subjectiv Betydning, og derfor maa Ordet
„\textit{intelligam}“ tages i selvsamme Betydning. „Jeg troer, for
at jeg kan erkjende“, efter Christi Forjættelse (Joh. 8, 31—
32): „Dersom I blive i mit Ord, ere I sandelig mine
Disciple, og I skulle erkjende Sandheden, og Sandheden
skal gjøre Eder frie.“ Denne Erkjendelse, der ikke objectivt
raisonnerer over Dogmet, men i inderlig og subjectiv Overbeviisning
tilegner sig det objective Dogme, kan alene stille
sig i et positivt Forhold til Dogmet. Men det maa vel
fastholdes og ikke forglemmes, at, idet den subjective Erkjendelse
stiller sig i sit positive Forhold til det bestemte Dogme,
stiller den sig med det Samme i et negativt Forhold til al
mulig objectiv videnskabelig Sandhedserkjendelse, af hvad
Navn nævnes kan.

% --- p. 36 ---
\setcounter{footnote}{0}%
\opage{36}Dette negative Forhold er nu deels et Ligegyldighedens
Forhold, forsaavidt den Troende praktisk holder sig til Dogmet,
uden at bekymre sig om Videnskabens objective Begreber;
deels et polemisk Forhold, forsaavidt den Troende tillige
bliver vidende om de Indsigelser, ved hvilke Videnskaben
kan forsøge at svække Dogmets religiøse Betydning. For
at erkjende Dogmet om Christi Sæde ved Faderens Høire,
tør den Erkjendende i Troens Navn aldeles ikke indlade sig
paa at ville udgrunde, hvorvidt en forklaret Individualitets
Allestedsnærværelse lader sig \emph{tænke} eller ikke \emph{tænke}. Troeserkjendelsen
bliver ved \textit{Ita}, og skrider ikke frem til \textit{Quare};
den bliver i Christi Ord, og indlader sig ikke paa nogen
Speculation.

Naar jeg da nu efter denne korte Forberedelse forsøger,
hvad Dr. M. fordrer, at „lade Erkjendelseslæren gaae Haand
i Haand med en Prøvelse af selve Dogmeudviklingen“, saa
spørges: hvilken Erkjendelseslære kan man da forefinde i Dr.
Ms Udvikling af Dogmet om Christi Sæde ved Faderens
Høire.

At Dr. M. i Ubiqvitetslæren vil corrigere de gamle
lutherske Dogmatikere, og at forsaavidt „det reformatoriske
Moment“ ogsaa her viser sig i hans Dogmatik, seer jeg
nok; men jeg seer ogsaa, hvad det er for et Correctiv, han
har anbragt. Man kan nemlig i det her omhandlede Punkt
corrigere de Gamle paa to Maader, alt eftersom man enten
beraaber sig paa Troeserkjendelsen eller paa den objective
videnskabelige Tænkning. I første Henseende viser man, at
de Dogmatikere, der fixerede Forestillingen om Christi Legemes
Allestedsnærværelse, just stode i Begreb med at overskride
den religiøse Erkjendelses Grændse, og kun forsaavidt
% --- p. 37 ---
\setcounter{footnote}{0}%
\opage{37}bleve frelste fra den objective Speculation, som de i Troens
Medfør lode Forstanden drage en Conseqvens, i hvilken
Forstanden selv gik tilgrunde. I sidste Henseende indrømmer
man derimod, at de Gamle vare paa rette Spor, idet de
gjorde „Christi Sæde ved Faderens Høire“ og hans individuelle
Nærværelse i den hellige Nadvere til Gjenstand for
objectiv Tænkning, men søger saa at rette paa deres objective
Begriben ved en endnu dybere gaaende Begriben.
Hvilken Vei Dr. M. er gaaet med sit „Reformatoriske“, kan
ikke drages i Tvivl: han er aabenbart gaaet den sidste.

At Dr. M., idet han „fremdeles“ vedbliver med sin
objective, men dog „paa religiøs Livserfaring“(!) grundede,
Erkjendelse af Aabenbaringens Sandheder“, fremdeles
paa det Bestemteste vil have den objective Erkjendelse i
Troen og udaf Troen adskilt fra en anden objectiv Erkjendelse,
der ikke er af Troen, har jeg heller aldrig overseet.
Der skal altsaa, ifølge Dr. Ms „Grundanskuelse“, gives tvende
Arter af objectiv Erkjendelse, nemlig: en lavere, verdslig, almeenvidenskabelig,
der endnu ikke kan sige: \textit{Credo, ut intelligam},
efterdi den endnu er i sin Syndighed, og i denne sin
Syndighed kun formaaer at føde den „kosmiske“, „pantheistiske“,
„rationalistiske“ Viden, for hvilken Troen er et Paradox; og en høiere, kirkelig, dogmatisk-videnskabelig Erkjendelse,
der strax kan sige: \textit{Credo, ut intelligam}, efterdi
den er forløst ved Guds Naade, og det saaledes forløst, at
den kan producere den ægte theologiske Speculation, for
hvilken Troen ikke er noget Paradox (at sige, naar man
tillige har saamange Sprogkundskaber, at man til Ordet
„Paradox“ kan søge det tilsvarende Udtryk i Modersmaalet).

Hvorledes Dr. Martensen vil have denne Adskillelse for%
% --- p. 38 ---
\setcounter{footnote}{0}%
\opage{38}staaet, udtaler han paa det Bestemteste i sine „dogmatiske
Oplysninger“, hvor han blandt Andet (S. 22) „henstiller“
til mig, „at overveie, om jeg fortrøster mig til at godtgjøre,
at Forløsningen af den faldne Menneskenatur ikke
ogsaa omfatter Tankens Forløsning, at dette ikke er Christendommens
Hensigt, fuldstændig at restituere den sande
Humanitet, men at lade væsentlige Elementer --- og Tanken
hører dog vel til de væsentlige Elementer --- forblive udenfor
sin helliggjørende og fuldkommende Virksomhed.“

I den speculative Prædikestiil, der prædiker Speculationen
uden at tænke Tanken, kan en saadan Forsikkring om
„Tankens Forløsning“ jo vistnok tage sig ret godt baade
opbyggeligt og speculativt ud; men --- „viis mig din Tro
af dine Gjerninger!“

Skulde Ideen om „Tankens Forløsning“, saaledes som
Dr. M. opfatter den, have nogen sand speculativ-dogmatisk
Gyldighed: da maatte jo denne Forløsning være for Dogmatikeren,
hvad Geniet er for Philosophen. Det maatte
da vise sig, at Dogmatikeren, hvis „Genius“ var bleven
„forløst“, virkelig formaaede, at behandle Aabenbaringens
Problemer med en saadan videnskabelig Dybsindighed og
udvikle Begreberne med en saadan Skarphed, at selv den
meest begavede Philosoph maatte misunde „den Forløste“
ikke blot hans „Evne til umiddelbar og levende Ideeanskuelse“,
men tillige hans vidunderlige Gave til høiere Tænkning
og indtrængende Begriben. Er dette nu i fjerneste
Maade anvendeligt paa den Martensenske Behandling af de
christelige Dogmer? Jeg henstiller til Læseren at overveie,
om han „fortrøster sig til at godtgjøre, at Dr. Ms videnskabelige
Udvikling af Ubiqvitetsdogmet bærer mindste Spor
% --- p. 39 ---
\setcounter{footnote}{0}%
\opage{39}af nogen overnaturlig „Tankens Forløsning.“ Thi jeg kan
virkelig ikke finde, at denne „dogmatiske Genius“ enten i §
177 eller i nogen anden dogmatisk Paragraph har truffet
ringeste videnskabelig Anstalt til at gjøre den paaberaabte
Adskillelse af „Christusnærværelsen“ og „Logosnærværelsen“
i nogen Grad indlysende for den objective Tænkning. Derimod
forekommer det mig, at Alt, hvad han i denne Anledning
„producerer“, er af den Beskaffenhed, at end ikke
den simpleste „Dilettant i Logiken skal misunde ham „Productiviteten.“


„Men“ --- indvender maaskee den „velvillige“ Læser ---
„hvad Du her siger, kan nu fra dit Synspunkt være
ganske rigtigt; men hvad beviser det? Det beviser jo kun,
hvad ogsaa Dr. Martensen (Dogm. Oplysn. S. 45) kan
bevidne, at Du fra Begyndelsen af har misforstaaet \textit{Credo},
\textit{ut intelligam}. Da Du nu saaledes, som Dr. Martensen
kan bevidne, fra Begyndelsen af saa letsindig misforstod
\textit{Credo, ut intelligam}, at Du uden at tænke over din Tro
--- \textit{quod negligentiæ mihi videtur}! --- endog sad og ventede
paa, at der „ved et Lykketræf skulde tilfalde Dig en
Arv af ontologiske og andre Erkjendelser“: hvad Under da,
at Du i din „Fremfusenhed“ og --- jeg maa sige Dig det
alvorligt, --- i din „Letsindighed“ ikke kan see det Betydningsfulde
i Dr. Martensens Productivitet: for at erkjende, hvad
den forløste Tænkning her producerer, maatte din egen Tanke
jo selv være forløst. Vi Andre derimod, der fra Barnsbeen
af ere opfødte og just derved gjenfødte i Christendommen,
vi Andre, der alt bleve vante til Paradoxet, da vi
bleve afvante fra Moders Bryst, alle vi Andre, der saadan
lidt efter lidt ere blevne fortrolige med Aabenbaringens
% --- p. 40 ---
\setcounter{footnote}{0}%
\opage{40}Sandheder, idet vi saadan lidt efter lidt have faaet „Tankens
Forløsning“, see, alle vi Andre, vi kunne vidne for Dr.
Martensen, at han, idetmindste relativt og til en vis Grad,
har begrebet baade Logosnærværelsen og Christusnærværelsen,
baade Christusnærværelsen ved Faderens Høire og Christusnærværelsen
i den hellige Nadvere, ja alle vi Andre
kunne for vor Deel vidne, at Dr. Martensen speculativt har
forstaaet Christendommen, som Dr. Martensen vistnok ogsaa
for sin Deel vil kunne vidne, at alle vi Andre speculativt
have forstaaet ham!“

Kjære Læser, hvo Du end er, gjenfødt eller ikke gjenfødt
(thi derom kan jeg i nærværende Tilfælde Intet vidne),
før Du end yderligere begeistrer Dig for Speculationens
\emph{vidnende} Methode: tillad mig da blot et lille Spørgsmaal:
Fortrøster Du Dig virkelig til at godtgjøre, at Dr. M. i
nogensomhelst dogmatisk Paragraph, og da navnlig (for at
dvæle ved dette enkelte Dogme) i Paragraphen om Christi
Allestedsnærværelse, har havt Mod til at overtage Conseqvensen
af sit eget speculative Vidnesbyrd?

Havde Dr. M. for Alvor selv troet paa den objective
videnskabelige Tankes „Forløsning“, saa maatte han jo have
behandlet „Christusnærværelsen“ paa en ganske anden Maade,
end skeet er. Han maatte da ikke blot have „vidnet“, at
Begrebet af Christusnærværelsen i Himlen og den individuelle
Christusnærværelse i Nadveren er det „naturlige Menneske
et Paradox“, men burde tillige have „vidnet“ udtrykkelig,
at han for sin Deel meget vel --- idetmindste til en
vis Grad --- kunde tænke dette overnaturlige Begreb speculativt
ved Hjælp „af Tankens Forløsning“; ja han burde
have vidnet, at baade Anselmus (den Mand, der tidlig
% --- p. 41 ---
\setcounter{footnote}{0}%
\opage{41}„indaandede ham Lyst til dogmatisk Speculation“) og „den
første Adam“ og „Peder paa Forklarelsens Bjerg“ virkelig
have kunnet tænke, hvad der ellers er den menneskelige Forstand
utænkeligt, Begrebet af Christi Legemes individuelle
Mangestedsnærværelse, og --- vel at mærke --- tænke dette
Begreb videnskabeligt, objectivt, contemplativt, speculativt,
dogmatisk. Kort sagt: Dr. Ms Paaberaabelse paa „Tankens
Forløsning“ indeholder den Paastand, at Christus og Apostlene
skulle have forjættet de Troende en ny speculativ Logik,
en objectiv Logik eller Ontologie, hvorved det siden hen,
naar de første evangeliske Mirakler ophørte, kunde blive muligt for en „dogmatisk Genius“ som „den hellige og almindelige
Kirkes Organ“ at gjøre et nyt og mere tidssvarende
Mirakel, et speculativt videnskabeligt Mirakel ved nemlig at
gjøre Det tænkeligt, som ellers enhver anden videnskabelig
Tænker maa finde utænkeligt.

Har nu Dr. M. virkelig, naar det kom til Stykket,
gjort Alvor af denne Paastand? Jeg kan ikke see det. Thi
det lille Tankeforsøg, den „dogmatiske Genius“ her anstiller
over Christi Legems Allestedsnærværelse, synes mig just ikke
at udvise nogen „helliggjørende og fuldkommende Virksomhed
eller bære Spor af nogen særegen Gjenfødelsens Logik.
--- „Selv en forklaret Individualitet, selv et aandeligt Legeme
lader sig ikke tænke uden Begrændsning.“ Denne
Tanke er virkelig saa beskeden, saa tilgjængelig, saa ligefrem,
naturlig og spagfærdig, at den ikke i mindste Maade overstiger,
hvad den almindelige Logik i saa Henseende kan
overkomme.

Naar nu Dr. M. med'alt Dette saa bitterlig klager over, at hans Recensenter ikke have „givet sig Tid til at lade
% --- p. 42 ---
\setcounter{footnote}{0}%
\opage{42}Bedømmelsen af Erkjendelseslæren gaae Haand i Haand med
en Prøvelse af selve Dogmeudviklingen, og vil, at Læseren
skal betragte Indledningens „almindelige Principer i uopløselig
Forbindelse med den virkelige Saganskuelse, idet han
paastaaer, at disse almindelige Principer først ved en saadan
Saganskuelse faae „Liv og Skikkelse“, da indseer jeg virkelig
ikke, hvad Fordeel den høiærværdige Forfatter kan
vente sig af at indskjærpe en saadan Fordring. Thi, jeg
vil ikke fordølge det, det forekommer mig, som om.
Dr. Ms „Erkjendelseslære“ tager sig langt bedre ud, naar
den holdes i en vis almindelig og svævende Afstand fra den
„virkelige Saganskuelse,“ end naar man anvender den til
„en Prøvelse af selve Dogmeudviklingen.“ Hvad Hs. H.
f. Ex. i de dogmatiske Oplysninger oplyser om „den umiddelbare
og levende Jdeeanskuelse,“ om „Tankens Forløsning“,
om „Forløsningsbevidstheden“ og „Aabenbaringsbevidstheden“,
om al „den Herlighed og Guddommelighed og
himmelske Klarhed,“ „Aabenbaringen“ formedelst „Aabenbaringsbevidstheden“
har i sig selv „paa Grund af det Indblik,
den lader os gjøre i Guds Dybheder“, hvilket Indblik
igjen stadfæstes ved speculative Henblik til Anselmus og den
første Adam og Peder paa Forklarelsens Bjerg, og det saaledes,
at disse Henblik igjen stadfæstes ved nye Indblik i
„den profetiske Anskuelse af den nye Himmel og den nye
Jord“, i „Legemlighedens Forklarelse“, i „Contemplationens
Rose“ og Speculationens „stille Glæde“ („I dit Lys skulle
vi see alt dette tykkes mig at være fremsat saa smukt,
saa dybsindigt, aandrigt og opløftende, at det næsten er
Synd at drage det ned fra det Svævendes ætheriske Høide
og sætte det i Forbindelse med nogen „virkelig Sagan%
% --- p. 43 ---
\setcounter{footnote}{0}%
\opage{43}skuelse“; thi ved Forbindelsen med Noget saa prosaisk som
en virkelig Saganskuelse vilde jo alle disse „Contemplationens
Roser“ falme og tabe deres Duft.

Hvordan det altsaa vil gaae med den Martensenske
„Erkjendelseslære,“ naar den virkelig skal „gaae Haand i
Haand med en Prøvelse af selve Dogmeudviklingen“, er nu
viist i Udviklingen af Ubiqvitetsdogmet. For at Ingen skal
mene, at jeg her vilkaarligen har valgt et Dogme, som
ifølge sin Natur hører til de meest ubegribelige, maa jeg
gjøre opmærksom paa, at jeg aldeles ikke vedkjender mig
den nylig opfriskede Adskillelse imellem „meer og mindre
begribelige“ Dogmer, og at Dr. M., hvis Skarpsindighed bør
have Æren for denne Distinction, selv maa indrømme, at
dette Dogme, ifølge hans egen Rangordning, hører til de
meest begribelige. „Mørkets Ubegribelighed“ --- siger han,
Dogm. Opl. S. 58 --- er væsentlig forskjellig fra Lysets,
hvoraf da ogsaa følger, at den Begribelighed, om hvilken
her kan være Tale, er en specifisk anden\dots{};“ men den forklarede
Christus tilhører jo ikke Mørkets, men Lysets Rige.
Dr. M. kan altsaa ikke beklage sig over, at jeg ved denne
Leilighed forvexler ham det „Teleologiske“ med det „Anti-Teleologiske;
thi „Christi kongelige Embede“ hører dog vel
med til „det Teleologiske.“ --- Hvad der imidlertid nærmest
har foranlediget mig til at dvæle saalænge ved Ubiqvitetsdogmet,
er ingenlunde den Martensenske „Modsætning imellem
det Teleologiske og det Anti-Teleologiske“, men, som sagt,
den Stillingske Anmærkning.

Idet jeg saa hermed slutter Forhandlingen om Mag.
Stillings Competance som Recensent af Dr. Martensens
Dogmatik, skal jeg endnu kun tillade mig en Bemærkning.

% --- p. 44 ---
\setcounter{footnote}{0}%
\opage{44}At Dr. M., naar han overhovedet følte sig oplagt til at
svare „sine Recensenter“, maatte svare i en skarp Tone, var
ifølge „det Sprog“, man havde tilladt sig mod Hs. H., at
vente. At Dr. M., i det Øieblik han nedlader sig til at betræde
Arena'en, strax viser sine Modstandere, at der er Riis
i den Kost, hvormed han nu agter „at feie for sin Dør“, er
polemisk i sin Orden.\footnote[1]{Saaledes er det f. Ex. ret passende og ingenlunde uvittigt, at han i sine dogmatiske Oplysninger præsenterer mig som Den, der vil arve de „rige Onkler.“ ---} At Dr. M. ikke lader det blive ved
det intellectuelle Riis, men (ret som om han uvilkaarlig
følte, at hans intellectuelle Riis ikke var saltet not) tillige
griber det moralske (sit \textit{venia verbo}) Kosteskaft, for at
virke med desto større „Alvor“, er vel ikke saa ganske i sin
Orden; dog, da Dr. M. nu een Gang har meent at burde
exseqvere den moralske Revselse, skal ogsaa denne Forret
for mit vedkommende være hans „dogmatiske Genius"
velvillig indrømmet, og Alt, hvad han ved denne Leilighed
har sagt om min „Paatrængenhed“, Fremfusenhed“, „Letsindighed“
og „Ustadighed“, for gammelt Venskabs og Bekjendtskabs
Skyld tilgivet.

Derimod er der een Ting, man paa det Alvorligste
maa foreholde Dr. Martensen, og det ikke blot for Dr. Martensens egen Skyld, men ogsaa for „Litteraturens“ Skyld
(thi Dr. Martensen og „Litteraturen“ og „Litteraturen“ og
Dr. Martensen kunne her ikke adskilles), og det er, at
Hs. H. ikke en anden Gang, naar han vil discutere
„de almindelige Principer“, forhaster sig
med at \emph{lægge} Modstanderens „Skriverie“ hen,
% --- p. 45 ---
\setcounter{footnote}{0}%
\opage{45}før han ret har seet sig for, hvad det var, der i
dette „Skriverie“ stod skrevet. At Dr. M. ved denne
Leilighed har meent at burde „holde Justits i Litteraturen“,
betvivler jeg ikke; men --- for en anden Gangs Skyld ---
var det dog virkeligt ønskeligt, om en Mand, til hvis Autoritet
Dr. Martensen sætter Lid, ret vilde indprænte ham
ovenstaaende velmeente Raad; thi hvis Hs. H. fremturer med
sit excommunicerende Magtsprog, som han her har begyndt,
kunde han jo let misbruge Banstraalen til stor Skade
baade for sig selv og for „Litteraturen.“

Som et Slags Undskyldning for Mag. Stillings Afviisning
kunde man maaskee betragte Dr. Ms Yttring
(Dogm. Opl. S. 8), at han „umulig kan tage Hensyn til
Alle, allerede af den Grund, at dette vilde føre til en trættende
og uoverkommelig Vidtløftighed. I den Forudsætning,
at Alt er afgjort, naar de „Sætninger“, paa hvilke min
„undersøgende Anmeldelse“ støtter sig, ere kuldkastede, mener
da den høiærværdige Forfatter at træffe Sømmet paa Hovedet,
ved at træffe mit Hoved, og altsaa er det da nu mig,
der alene maa holde for.

Uden at ville opholde sig ved min „Skrivemaade“ (hvad
der vel ogsaa maa ansees for overflødigt, efterat H. H. selv
har givet „Beviis paa sin Skrivemaade“) vil Dr. M. blot
opholde sig lidt over „Form og Tone“ i min Anmeldelse,
og gaaer saa uden Omsvøb løs paa Sagen, idet han erklærer:

„De Sætninger, Recensenten (Prof. N.) opstiller imod
mig, betragter jeg som Sætninger, for hvilke Ingen uden
han selv er ansvarlig. Om hine Sætninger maaskee have
en heelt anden Mening i den humoristiske Sammenhæng
% --- p. 46 ---
\setcounter{footnote}{0}%
\opage{46}hos Pseudonymerne, om Recensenten her er i Forstaaelse,
eller i Ikke-Forstaaelse, eller i Misforstaaelse, alt Dette ligger
uden for Omkredsen af min Undersøgelse. Jeg skal her kun
indstrænke mig til den langt ringere og beskednere Opgave
at feie for min egen Dør ved at vise, at jeg aldeles ingen
Brug har for de Læresætninger, Recensenten med saamegen
Paatrængenhed har villet meddele mig, efterdi jeg paa disse
maa anvende det gamle Ord, at det Sande deri ikke er
Nyt, og det Nye ikke Sandt; ligesom jeg ogsaa skal vise, at
hvorledes det end monne forholde sig med hans Forstaaelse
af Pseudonymerne, saa er det uomstødeligt vist, at han ikke
har forstaaet mig, men paa det vilkaarligste misforstaaet og
mistydet mine Tanker“ (Dogm. Oplysn. S. 13 ---14).

At min „undersøgende Anmeldelse“ i visse Puncter er
noget nærgaaende, vil jeg ikke negte; men den „velvillige“
Læser vil da vist heller ikke negte, at Anmeldelsens blot
sammenstillende, undersøgende og spørgende Form er lemfældig
og svag i Forhold til de stærke Ord, i hvilke
Dr. Martensen her proclamerer sin dogmatiske Speculations
Uantastelighed.

Anmeldelsen, der begynder med „en Undskyldning, at
den tillader sig at sammenstille tvende, i Retning, Aand og
Tone saa ueensartede Skrifter,“ som Mag. S. Kierkegaards
„Johannes Climacus“ og Dr. H. Martensens „Christelige
Dogmatik,“ deler sig i to Dele.

Anmeldelsens første Deel (S. 9—12) indeholder kun,
hvad Dr. M. kalder „et Excerpt“ af Pseudonymen, en sammentrængt
Fremstilling af de Hovedtanker, der bevæge Humoristen,
idet han tager sig den Frihed at behandle Theologernes
Forsøg med den objective (historiske og speculative)
% --- p. 47 ---
\setcounter{footnote}{0}%
\opage{47}Viden af Christendommens objective Sandhed, ikke som en
videnskabelig Opgave, der egentlig bør discuteres i Videnskaben,
men som den Art „alvorlige“ Bestræbelser, der, naar
de sees i Ideen, vise sig at være af den Beskaffenhed, at
de væsentlig høre hjemme i den høiere Komik. Hvad den
komiske Udførelse selv angaaer, kunde denne naturligviis ikke
gjengives i et „Excerpt“; hvorfor jeg da heller ikke har tilladt
mig noget Forsøg i denne Retning. Men, at der igjennem
„Humoristens“. Komik gaaer en frygtelig Alvor, var
mig tillige saa klart, at jeg, især da jeg saae, hvorledes
Dr. Martensen i „Forordet“ til sin Dogmatik var kommen
paa og talte saa betydningsfuldt om „Strøtanker
og Aphorismer, Indfald og Glimt“ --- umulig kunde undlade
at henlede Hs. Høiærvs. Opmærksomhed paa det
Spørgsmaal: \emph{Har ikke „Johannes Climacus“ dialektisk
skjærpet det christelige Problem?}

Da nu saaledes Anmeldelsens første Deel udelukkende
beskjæftiger sig med Pseudonymens Theses, og ikke med et
eneste Ord falder ind i den Martensenske Tankegang, kan
Beviset for, at jeg har misforstaaet Dr. M., ikke hentes fra
dette Parti. Hele Eftertrykket falder altsaa paa Anmeldelsens
anden Deel, der (S. 42 --- 132) heelt igjennem bevæger
sig om det Spørgsmaal:

\emph{Har ikke „den christelige Dogmatik“ udialektisk
omgaaet det christelige Problem?}

I denne anden og forholdsviis udførligere Afdeling (der
endog ved en liden Streg er adskilt fra den første) har jeg
da --- saavidt jeg veed --- heller ikke brugt nogen fremmed
Tunge, men, med Undtagelse af Citater, som ikke vel tillade
% --- p. 48 ---
\setcounter{footnote}{0}%
\opage{48}nogen Omskrivning, fremsat Alt med mine egne Ord, og
forklaret Dr. M. min Mening om hans Dogmatik, saa godt
jeg har kunnet.

Vil Dr. M. altsaa gjøre det „uomstødeligt vist“, at jeg
ikke har forstaaet ham, men „paa det Vilkaarligste misforstaaet
og mistydet hans Tanker“: da er jeg vel ogsaa berettiget
til at vente, at Hs. Høiærv. deels indlader sig paa
de i min Recension udviklede Hovedpuncter, deels begrunder
sit eget „mistydede“ Videnskabsbegreb saa tydelig og bestemt,
at deraf kan sees, at min Opfatning har været en Misforstaaelse.
Maaskee vil Dr. M. engang med Tiden gjøre begge
Dele; men, at han endnu i sine „dogmatiske Oplysninger“
hverken har gjort det Første eller det Sidste, er --- uomstødeligt
vist.
\emph{Dr. Martensen har jo i sine dogmatiske Oplysninger aldeles ikke indladt sig paa at besvare}
de Indvendinger, min undersøgende Anmeldelse
fremsætter mod hans Dogmatik.

For strax at angive Forskjellen imellem Troeserkjendelsen
og den objective Viden, begynder Anmeldelsens anden
Deel med en Undersøgelse af, hvad det vil sige, at være
„absolut i Troen“ og „absolut i Viden.“ Forfatteren af de
dogm. Opl. er imidlertid saa langt fra at tage noget Hensyn
til denne Undersøgelse, at han endog tillægger mig saadanne
Meninger, som han umulig kunde have tillagt mig,
hvis han i mindste Maade havde ændset mine Ord.

Hvad t. Ex. Viden angaaer, hedder det udtrykkelig % PRINTED AS IS: t. Ex. (for f. Ex.)
(Anm. S. 49): „At erklære sig absolut i Viden er en Anmasselse,
forsaavidt den Speculerende dermed vil sige sig at
have en i alle Henseender absolut og fuldkommen Viden af
% --- p. 49 ---
\setcounter{footnote}{0}%
\opage{49}det Absolute\dots{} Ikke destomindre lader Dr. M., som om
Sligt aldrig var sagt, og paadutter mig uden videre den
Anmasselse, at ville gaae ud fra en Speculation, der paatager
sig at gjøre Guds Væsen absolut begribeligt. --- „Men,
svarer Prof. N., .... Gud maa enten være absolut begribelig
eller absolut ubegribelig.“ (Dogm. Opl. S. 52).

Hvad Troen angaaer, siger Anmeldelsen (S. 50): „At
ville gaae og gjælde for absolut i Troen er en Anmasselse,
hvis den Enkelte dermed vil udgive sig for at kunne fuldkomme
al Troens Gjerning, at have en i alle Maader absolut
og ubetinget Tro, den Tro, som flytter Bjerge“
Men, langt fra at gaae ind paa den i Anmeldelsen fremhævede
Adskillelse af, at holde sig absolut til Troens Princip,
og, at udgive sig selv for en absolut Troende (en Adskillelse,
som Forordet til Dogmatiken gjorde høilig fornøden),
gaaer Forf. af dogm. Opl. dette Punkt med tilsyneladende
Uvidenhed forbi, og tillægger mig uden videre den Anmasselse,
at ville udgive mig for Troeshelt (jfr. dogm. Opl. S.
71—72 o. o. St.), og dette har Hs. H., som det synes,
gjort, ene og alene, for at kunne (Opl. S. 12) have den
Fornøielse at gjentage den Tvetydighed, han har anbragt i
Forordet til sin Dogmatik, „at ligesom der var en Tid, da
vi havde ikke Faa iblandt os, som altfor tidligt vare blevne
absolute i Viden, saaledes fattes det nu ikke paa Saadanne
iblandt os, som altfor tidligt ere blevne absolute i Troen\footnote[1]{At en Hr. Cand. Hjort kan komme springende med den Nyhed, at jeg nu virkelig bilder mig ind at være den Troende selv, fordi jeg i mit Skrift „Evangelietroen og den moderne Bevidsthed“ har ladet den Troende tale, faaer jeg naturligviis at finde mig i; men, at Dr. Martensen nu ogsaa tager denne Opdagelse til Indtægt, falder mig lidt tungt. Dr. M. burde dog virkelig kunne gjøre Forskjel mellem det, at lade den Troende tale, og det, at tale som den Troende; Dr. M. siger jo desuden selv, at han har læst det Skrift af mig, hvis Forord ender saaledes: „Forfatteren af nærværende Skrift tilstaaer oprigtig, at han endnu hverken i Ord eller Gjerning ret forstaaer at stride for Troen, som det sig bør“ o. s. v.}
og altfor tidligt ere blevne Troens Forsvarere.“
% --- p. 50 ---
\setcounter{footnote}{0}%
\opage{50}Som Bidrag til nærmere Bestemmelse af Forholdet
mellem Troeserkjendelsen og Videnskaben behandler Anmeldelsen under den Overskrift:

„Den christelige Dogmatik“ famler efter Problemet, men

% CHECK p. 50: centred line: heading or verse
\begin{center}finder det ikke,\end{center}

det Spørgsmaal, hvorvidt den speculative Dogmatik væsentlig
lader sig adskille fra den saakaldte Religionsphilosophie
eller ikke. Besvarelsen af dette Spørgsmaal er i nærværende
Sag af gjennemgribende Betydning. Lod det sig virkelig
paavise, at der i Henseende til Princip, Methode og videnskabelig
Interesse bestod en sand og betydningsfuld Modsætning
imellem den speculative Dogmatikers og den speculative
Philosophs Opfattelse og Behandling af de christelige
Dogmer: saa fik Dr Martensen i Sandhed Ret, saa vilde
man virkelig dog tilsidst blive nødt til at indrømme, at der
gaves en dobbelt Speculation, en dogmatisk og en philosophisk,
og at den første i sin Art var ligesaa videnskabelig
berettiget, som den sidste. Nu stræber „Anmeldelsen“ at godtgjøre,
at en saadan Modsætning kun bestaaer i Dogmatikerens
Indbildning, ikke i Sagen. „Anmeldelsen“ efterviser,
at det Forsøg, Dr. M. har gjort med at adskille Dogmatikens
Begreb fra Religionsphilosophiens, opløser sig --- trods
alle hans Forsikkringer og Paastande --- i et reent Intet;
ja det fremhæves udtrykkelig, at den Martensenske Adskillelse af „id%
% --- p. 51 ---
\setcounter{footnote}{0}%
\opage{51}eologisk“ og „philosophisk Speculation“ har havt et
saa uheldigt Udfald, at endog tvende Paragrapher i hans
egen Dogmatik desangaaende ere geraadede i aabenbart
Fjendskab med hinanden,

Hvad har nu Forf. af dogm. Oplys, oplyst i denne
Anledning? Intet, aldeles Intet. Uden at ændse de mod
hans Adskillelse fremsatte Indvendinger, gjentager han blot
sin Paastand, at den theologiske Speculation er af en ganske
anden og høiere Art end den philosophiske, paanøder eller
rettere paadigter mig (Opl. S. 14) den Antagelse, at der
ved den objective Viden blot skulde tænkes paa „en ligegyldig
Viden“, uagtet „Anmeldelsen“\footnote[1]{„Men, naar det nu er vitterligt, at denne Adskillelse (imellem Dogmatikerens og Philosophens Interesse for Erkjendelsen af Dogmet) beroer paa en ulyksalig Forvexling af det Personlige og det Geniale; naar det nu er vitterligt, at,En kan være en troende Christen og arbeide kraftig i Troen, uden at have den intellectuelle d. v. s. den theoretiske Kjærlighed, hvorom her tales, medens en Anden kan være aldeles behersket og betagen af denne Kjærlighed, uden i samme Grad at arbeide for den subjective Tilegnelse i Troen; naar det er vitterligt,\dots{}. at den intellectuelle og speculative Kjærlighed, hvormed Tænkeren fordyber sig i Aabenbaringens Mysterier, ingenlunde er nogen egentlig religiøs Fordring („Troen er jo det ene Fornødne for os Alle“), men ene og alene Talentets Drift og Trang til Beskuelse af Ideen, til objectiv Glæde i guddommelige Syner: hvor er saa Forskjellen?\hfill Unders. Anm. S. 65.} siger lige det Modsatte; % PRINTED AS IS: at,En, in the note
ja gaaer endog (Oplys. S. 27) strengt i Rette med
mig, fordi jeg skal have overseet, at enhver Viden maa
have sin Interesse, da „en absolut interesseløs Viden neppe
lader sig tænke“. --- Saaledes forstaaer nu Dr. M. sine
„Recensenter“. Man gjør ham opmærksom paa Forskjellen
imellem den religiøse og den videnskabelige Interesse, deraf
udleder han, at man taler om „en absolut interesseløs
Vi%
% --- p. 52 ---
\setcounter{footnote}{0}%
\opage{52}den“, og klager saa over, at man „misforstaaer og mistyder
hans Tanker“.

„Den christelige Dogmatik vil tilegne sig Speculationen,
men uden at fatte Speculationens Problem“.

Under denne Overskrift gaaer „Anmeldelsen“ (S. 67—90)
ind paa en nærmere Udvikling af den principielle Forskjel
mellem Troeserkjendelsen og Speculationen, der i det Foregaaende er paapeget. Undersøgelsen anlægger ikke her nogen
fremmed Maalestok, men finder overalt sit Udgangspunkt
Dr. Ms egen Dogmatik, det er Dr. Ms egne „Grundsætninger“
om Erkjendelsen i Troen og udaf Troen, der, under
omhyggelig Sammenstilling af herhen hørende Citater,
gjøres til Gjenstand for videnskabelig Prøvelse. Efter at
have sikkret sig Begrebet af, hvad Dr. M. kalder „den rene
Videnskab“, samler „Anmeldelsen“ sin hele Opmærksomhed
paa det Sted i „Dogmatiken“, der, hvad Standpunktet angaaer,
vel med Rette maa betegnes som det egentlige Hovedsted.


„For Dogmatiken er Christendommens absolute Sandhed
given forud og uafhængig af al Speculation. Hiint
δὸς ποῦ στῶ, som den søgende Philosophie saa ofte har
udtalt er oprindelig besvaret for Dogmatiken; og ikke gjør
Dogmatikeren Christendommen afhængig af sin Granskning,
men søger kun igjennem Tænkningen en dybere Tilegnelse
af den Sandhed, der er ham det Visseste af Alt, og til
hvilken han er kommen ad en heelt anden Vei end Speculationens.“
(Dogm. S. 5.)

Den skjulte Modsigelse, der ligger i disse Ord, drager
Anm. frem, og concentrerer sit Resultat i følgende Udtryk
(S. 80):\dots{} „er Troens Tænkning forskjellig fra Specula%
% --- p. 53 ---
\setcounter{footnote}{0}%
\opage{53}tionens, eller ere de i Grunden Eet? Indrømmer Dogmatikeren
det Sidste, saa tilstaae vi vel, at han har Ret,
naar han søger Troeserkjendelsens Fuldendelse i Speculationen;
men maae vi saa spørge, hvad der berettiger ham
til den Paastand, at han selv er kommen til den høieste
Sandhed ad en heelt anden Vei end Speculationens? Erklærer
han sig derimod for den første Antagelse, at nemlig
Troeserkjendelsen i Princip og Væsen er aldeles forskjellig
fra den speculative Erkjendelse, da indrømme vi ham gjerne,
at han forud maa være kommen til den absolute Sandhed
ad en heelt anden Vei end Speculationens; men maae vi
saa spørge, hvor han som Dogmatiker da vil hen med Speculationen?“


Hvad har nu Forfatteren af dogm. Opl. gjort for at
belyse dette Punkt? Han har (Opl. S. 51) paadigtet mig
følgende Feilgreb:

„Hvad da angaaer Prof. Ns Paastand, at jeg skal
have misforstaaet Speculationens Problem, da beroer denne
aabenbart derpaa, at han opstiller et heelt andet Begreb om
Speculationen, end det, jeg har vedkjendt mig i min Dogmatik,
og derved fra først af forrykker det hele Synspunkt.“

Altsaa: jeg holder mig til Dr. Ms eget i Dogmatiken
„opstillede“ Speculationsbegreb, for at vise, at dette Speculationsbegreb
i enhver Henseende er falsk og uholdbart; uden
at indlade sig paa mine Grunde, erklærer Dr. M., at jeg
opstiller „et heelt andet Begreb om Speculationen, end det,
han har vedkjendt sig“, og vil saa, at denne Erklæring skal
gjøre det „uomstødeligt vist“, at jeg „ikke har forstaaet ham,
men paa det Vilkaarligste mistydet og misforstaaet hans
Tanker.“

% --- p. 54 ---
\setcounter{footnote}{0}%
\opage{54}Ja Forfatteren af dogm. Opl. gaaer endnu lidt videre.
Han vedbliver nemlig at reise Beskyldninger, der ere af den
Beskaffenhed, at jeg næsten maa antage, at H. H. enten har
læst min Anmeldelse, som han, efter egen Tilstaaelse, har
læst Pseudonymerne, d. v. s. blot flygtig og „fragmentarisk“,
eller ogsaa glemt igjen, hvad han havde læst.

I Anledning af, at jeg (hvad, som ovenfor er viist, er
en usandfærdig Paastand) skulde være gaaet ud fra en Speculation,
der paatager sig at gjøre „de guddommelige Ting
absolut begribelige“, tillægger Dr. M. mig saaledes (Opl.
S. 52) følgende Argumentation: „Gud maa enten være
absolut begribelig eller absolut ubegribelig. Thi Kant og
Hegel vare systematiske Hoveder, og den ene af disse lærte,
at de guddommelige Ting vare absolut ubegribelige, den
anden, at de vare absolut begribelige. Ville vi derfor være
systematiske Hoveder, da maae vi enten lære det Ene eller
det Andet.“

Hvad jeg i min Anmeldelse (S. 124) har sagt, er
imidlertid kun, hvad der længst er anerkjendt i Videnskaben,
at Den, der vil skrive et System og gjælde for et systematisk
Hoved, maa fastholde et Princip, og fra Først til Sidst
respectere dets Conseqvenser, hvad enten han saa forøvrigt
er Theolog eller Philosoph. For at betegne, hvad jeg mener
med det Udtryk: „et systematisk Hoved“, har jeg exempelviis
anført Kant og Hegel. Heraf tager Dr. M. Anledning
til at irettesætte mig for en \textit{metabasis in aliud genus}.

„Men“ --- hedder det (Opl. S. 52) --- „istedetfor at
beraabe sig paa Kant og Hegel, hvis Paaberaabelse i nærnærende % PRINTED AS IS: nærnærende (nær- doubled over the line end)
Sammenhæng er en \textit{metabasis in aliud genus},
% --- p. 55 ---
\setcounter{footnote}{0}%
\opage{55}skulde han (Prof. N.) hellere have fremført mig Dogmatikeres
Exempel af tilsvarende Auctoritet.“ ---

Den Læser, som har mindste Kjendskab til den nyere
Tids theologiske Litteratur, vil uden Tvivl studse og falde
i Forundring over den besynderlige Dristighed, hvormed Dr.
M. her taler om \textit{metabasis in aliud genus}. Det er jo
dog vel vitterligt, ja --- „uomstødeligt vist“, at Dogmatikerne
Daub og Marheineke væsentlig holdt sig til samme
\textit{genus cognoscendi}, som Philosopherne Schelling og Hegel,
og at Dogmatikeren Schleiermacher paa sin Viis var ligesaa
systematisk-subjectiv, som Philosophen Kant.

Det maa endydermere vække Forundring, at Forf. af
dogm. Opl. saaledes kan dadle mig som den, der ikke har
„fremført ham Dogmatikeres Exempel af tilsvarende Auctoritet“,
da jeg dog i min Anmeldelse (S. 89) udtrykkelig har
anført ham Daub og Marheineke, og, S. 66, mindet ham
om Schleiermacher.

Ikke mindre paafaldende forekommer mig den Sufficance,
hvorved Dr. M., uden i fjerneste Henseende at ændse,
hvad der staaer i Anmeldelsen (S. 82---84), fremsætter den
Fordring, at jeg „i Gjerningen skulde have gjendrevet den
Sætning, han har opstillet i sin Dogmatik, at ligesom de
altid ere blevne tilskamme, der have givet sig ud for at
have begrebet Alt, saaledes ere de ikke mindre blevne tilskamme,
der have meent een Gang for alle at kunne afstikke
Grændserne for al menneskelig Begriben; thi bag efter
viste det sig bestandig, at der gaves \textit{plus ultra}, og de formeentlig
faste Grændser aabenbarede deres bevægelige Natur
ved at flytte sig“ (Dogm. § 33).

„Ja“ --- svarer Anmeldelsen (S. 82) --- „ved at flytte
% --- p. 56 ---
\setcounter{footnote}{0}%
\opage{56}sig!“ og tilføier derpaa Adskilligt om den „flydende Grændse“
og om „\textit{plus ultra}“.

Dersom det at „finde sig selv“ og blive sit Standpunkt
tro skal vise sig deri, at man, hvad der saa end indvendes
imod den een Gang opstillede „Grundanskuelse“, blot holder
Ørene stive og trodser alle Indvendinger ved en dristig
Gjentagen af egne Paastande: da indrømmer jeg gjerne, at
Dr. Ms Adfærd i denne Sag kan kaldes dogmatisk-normal,
og at Hs. H. i saa Henseende maa siges at have
„fundet sig selv“. Skal man derimod lægge mindste Vægt
paa Sandhed og Sandhedsundersøgelse: da indseer jeg virkelig
ikke, hvad Dr. M. kan mene at opnaae ved en saa tillidsfuld
Gjentagelse af Paastande, hvis Uholdbarhed er saa
iøinefaldende, at de end ikke finde Tilhold i den populære
Mening.

At de altid „ere blevne tilskamme“, der have meent
een Gang for alle at kunne „afstikke Grændserne for al
menneskelig Begriben“, er en ligesaa ubestridelig som triviel
Sætning, naar Talen er om saadanne Ting, som virkelig
kunne være Gjenstand for „menneskelig Begriben“. Naturvidenskaben,
Lægevidenskaben, Mathemathiken, Philosophien, % PRINTED AS IS: Mathemathiken
kort, alle virkelige Videnskaber befinde sig jo i en stadig
Fremadskriden, i alle virkelige Videnskaber gives der \textit{plus}
\textit{ultra}, og Den maatte være forrykt, som her vilde paatage
sig een Gang for alle at „afstikke Grændsen for al menneskelig
Begriben“.

Men hvo seer nu ikke, at her i denne Sag tales om
at „afstikke“ en ganske anden „Grændse“. Her tales jo om
at „afstikke“ den Grændse, hvorved den indbildte Viden
nødvendigviis adskiller sig fra den sande Viden. Denne
% --- p. 57 ---
\setcounter{footnote}{0}%
\opage{57}Grændse er saare let at afstikke: Om det Overnaturliges
Væsen gives der kun en indbildt Viden; om den Lov, hvorefter
det Overnaturlige virker ind paa det Naturlige, gives
kun en indbildt Viden; om Miraklets objective Gyldighed
gives der kun en indbildt Viden; om den aabenbarede Religions
objective Sandhed gives der kun en indbildt Viden.

At den her „afstukne Grændse“ staaer fast, godtgjøres
saavel af Sagens Natur, som af Historiens Vidnesbyrd.
Det ligger i Videnskabens Væsen, at den ikke kan vove sig
ind paa det Overnaturliges Enemærker uden at give sig
Phantasien og Phantasteriet i Vold. Det sees af Theologiens
Historie, at denne Videnskab, langt fra at gaae fremad
i objectiv metaphysisk Viden, meget mere gaaer tilbage i
samme Forhold, som de virkelige Videnskaber gaae frem.
Naturvidenskaben, Philosophien og den historiske Kritik have
jo netop i den nyere Tid, hver paa sin Maade, bragt Theologien
i en saadan Selvmodsigelse, at der just ikke behøves
stor Skarpsindighed for at see, at Theologernes middelalderlige\footnote[1]{Det er i denne Henseende da heller ikke uden Betydning, at Dr. M., for at finde et Slags „Forbillede“ for sin „dogmatiske Speculation“, gaaer tilbage til Anselmus og Middelalderens Klostre.}
Videnskabsbegreb opløses Dag for Dag, og viser
sig lige ubrugeligt for Videnskaben og for --- Kirken.

„Den christelige Dogmatik vil tilegne sig Troen, men
uden at fastholde Troens Problem.“

For at retfærdiggjøre denne Paastand, slutter „Anmeldelsen“
sig (S. 90—125) saa omhyggelig og nøiagtig til
de Grundsætninger“, den Martensenske Dogmatik „opstiller“
om Troen, at man af det hele dogmatiske Værk ikke skal
% --- p. 58 ---
\setcounter{footnote}{0}%
\opage{58}kunne anføre mig et eneste Sted, hvis Betydning i dette
Punkt er overseet. Det er overhovedet ret charakteristisk,
at den Martensenske Dogmatik, hvis „Speculation“ jo er i
Troen, gaaer ud fra Troen og ud paa Troen, netop forraader,
hvad Magister Stilling i sin Recension (S. 29)
træffende bemærker --- „en paafaldende Mangel paa qvalitativt
betegnende Prædikater for Tro og Haab.“ Dette,
siger jeg, er charakteristisk; thi, var Dr. M. virkelig bleven
opmærksom-paa Det, der egentlig udgjør Troens Qvalitet:
han havde nødvendigviis maattet opgive „Speculationen.“

At Troen i sig forener baade det Subjective og det
Objective, at den Troende maa holde sig til Troens Gjenstand,
maa vide, hvad det er, han troer, og gjøre sig Rede
for sin Troes Indhold, forstaaer sig af sig selv, og er heller
ikke blevet omtvistet. Men, just idet man fra begge Sider
indrømmer, at Troen og Troens Gjenstand maa svare til
hinanden, saa gjentager sig her det Spørgsmaal, som i
nærværende Sag er det egentlige Hovedspørgsmaal: Hvorledes
skal Forholdet imellem Troen og Troens Gjenstand
rettelig bestemmes?

I Dr. Ms Dogmatik bestemmes Forholdet paa een
Gang som et relativt Vidensforhold og derhos tillige som
et Existentialforhold, d. e. et subjectivt personligt Forhold.

At denne Opfattelse af Forholdet er uklar, tvetydig og
selvmodsigende, stræbe Anmeldelsen at vise, ikke ved at anlægge
en fremmed og vilkaarl=Maalestok, men ved at holde % PRINTED AS IS: a damaged sort, „=“ overprinted by a curl, where vilkaarlig is meant
sig til Dogmatikens egne Ord.

„Dogmatikens anden Hovedsætning“ --- hedder det, Anm.
S. 96 —„gaaer ud paa at forebygge „den Misviisning, der
af Bekymring for „den Troende“ og hans Sjælstilstand bliver
% --- p. 59 ---
\setcounter{footnote}{0}%
\opage{59}ligegyldig for Troen (\textit{fides, quæ creditur}), der af Iver for
det Opbyggelige glemmer Indholdet, hvormed der skal bygges,
glemmer Grundvolden, hvorpaa der skal bygges.“ --- „Hvad
Reformationen vilde“ --- hedder det (Dogm. S. 40 § 21) ---
„var hverken udelukkende det Objective eller det Subjective,
det var den frie Forening af det Objective og det Subjective,
af Troens Indhold og Troens Inderlighed, af guddommelig
Aabenbaring og religiøs Selvbevidsthed.“

Imod denne „frie Forening af det Subjective og det Objective“,
saaledes som Dr. M. har forstaaet den, gjør Anmeldelsen
den dobbelte Indsigelse, at nemlig a) „det Objective“ herved
opløser sig, saa at det, der bliver tilbage, kun er et vagt og
ubestemt subjectivt Noget, og b) at „det Subjective“ fra sin
Side ligeledes forflygtiges saaledes, at det, der bliver tilbage,
kun er et ubestemt og ubestemmeligt objectivt Noget.

a) Hvad det første Punkt angaaer, vil Dr. M. naturligviis
sikkre sig „det Objective“ ved en videnskabelig Paaberaabelse
paa „Kirkens, Skriftens og Aandens Vidnesbyrd.“
Anmeldelsen bebreider ingenlunde Dr. M., at han
i sin Dogmatik har henviist os til disse Vidnesbyrd; men
den gjør ham opmærksom paa, at han ikke har forstaaet at
afbenytte disse saaledes, at enten Troen selv (\textit{fides, qua creditur}) eller Troens Indhold (Lærebegrebet --- \textit{fides, quæ}
\textit{creditur}) derved i mindste Maade klares og sikkres. Anmeldelsens
Tankegang er i Korthed følgende:

„Den objective Canon for al Christendom“ (Dogm. §
22) er Christus selv. Men Christus er kun i sin Kirke:
altsaa er det Kirken, „hin store Enkelte!“ --- der er den
objective Canon. Men Kirkens egentlige „Prøvesteen (lapis
lydius)“ er kun i Skriften: altsaa er det Skriften, der er
% --- p. 60 ---
\setcounter{footnote}{0}%
\opage{60}den objective Canon. Men Skriftens rette Forstaaelse er
kun i Aandens Vidnesbyrd: altsaa er det Aandens Vidnesbyrd,
der er den objective --- nei det er ikke sandt ---
Aandens Vidnesbyrd er en absolut subjectiv og „indvortes
Canon.“

Alt, hvad her siges om „det Objective“, gaaer altsaa
tilbage i det aldeles Subjective; thi Aandens Vidnesbyrd
er i det Subjective, ikke tildeels i det Subjective og tildeels
i det Objective, men ene og alene i det Subjective, absolut
i det Subjective. Hvor let er det ikke at see dette, naar
man kun vil see. Alle Theologer, Rationalisterne og de
Speculative ligesaavel som de Orthodoxe, ere jo i Kirken,
ville jo være i Kirken, beraabe sig paa Kirken og forsvare
Kirken, og det uagtet deres Begreber saavel om Christus,
som om Skriften og Kirken, ere saa modsigende og modstridende,
som tænkes kan. Skal nu Kirken, --- „hiin store
Enkelte“! virkelig antages at omslutte, vedkjende og tilegne
sig alle disse indbyrdes hinanden forkjættrende og fordømmende
Retninger: da er Aandens Vidnesbyrd i „den hellige
og almindelige Kirke“ tilvisse ikke blot et alsidigt, men tillige
et saa tvetydigt og tvetunget Vidnesbyrd, at Den, der
rettelig skal udtyde dette Vidnesbyrd, maa jo have et endnu
høiere „Aandens Vidnesbyrd.“

„Og hvad“ --- spørger Dr. M. (Dogm. S. 57) ---
„var Rationalismen Andet, end en stor \textit{theologia}
\textit{irregenitorum}, der oversvømmede den protestantiske Christenhed?“
Det objective Spørgsmaal om Forskjellen mellem den sande
og falske Theologie, der i Dogmatiken skulde afgjøres videnskabeligt,
forvandles altsaa dog umærkeligt til et subjectivt
og uvidenskabeligt Spørgsmaal om: hvilke Theologer der ere
% --- p. 61 ---
\setcounter{footnote}{0}%
\opage{61}„gjenfødte“, og hvilke der ikke ere gjenfødte. See, her have
vi det vage, ubestemte subjective Noget, der dog tilsidst skal
gjøre Udslaget, naar man har taget Troesobjectet forfængeligt ved at behandle det som et andet ligefrem videnskabeligt
Object, og just derved tabt alle objective Skillemærker.

Hvad har nu Forfatteren af dogm. Oplysn. gjort for
at belyse dette Punkt? Fremsat (Oplysn. S. 16) den sindrige
Bemærkning, at „selv den fattigste rationalistiske Dogmatik“
maa erkjende, at Sandheden har en vis Objectivitet“..,
og derhos --- skjøndt uden at anføre noget Sted,
der er forbigaaet --- (Dogm. Oplysn. S. 41) henkastet den
Beskyldning, at man „aldeles ikke gaaer ind paa Modstanderens
Tankegang“, ikke undersøger „Troesbegrebets Genesis
hos den Forfatter, mod hvem man stiler sit Angreb“, men
„uden videre antager et andet Troesbegreb som det sande,
omgaaes dermed som en fast Forudsætning, der er hævet
over al Tvivl“ o. s. v.

b) Betræffende det andet Punkt, overseer Anmeldelsen
ingenlunde, at Dr. M. ogsaa vil hævde „det Subjective“,
at han kalder Troeserkjendelsen en existentiel Erkjendelse, og
vil, at „det religiøse Gudsforhold skal betragtes som et Existentialforhold;
men --- \textit{hinc illæ lacrymæ}! just fordi Dr. M.
vil beraabe sig paa den existentielle Tænkning, just derfor bebreider
„Kritiken“ ham, at han befatter sig med at klare
denne Tænknings Indhold i Speculationens objective Viden;
thi den speculative Tænkning er i Princip og Væsen lige det
Modsatte af den existentielle. For den existentielle Tænkning
er det Subjective, det Personlige, det absolut afgjørende; i
den speculative Tænkning gaaer Tænkeren ud af sig selv,
% --- p. 62 ---
\setcounter{footnote}{0}%
\opage{62}taber sig i sin Gjenstand og forholder sig til den saa objectivt
og upersonligt, som det vel er muligt.

„Naar Dogmatikeren“ --- siger Anmeldelsen, S. 101, ---
„kommer ind paa det subjective Spor, anerkjender og indskjærper
han Subjectiviteten paa det Bedste, naar han saa
igjen falder over i den speculative Tankegang, hævder han
Objectiviteten, den objective Viden, „den theologiske Speculation“
paa det Eftertrykkeligste, og dog ere disse to Retninger
ikke blot hinanden aldeles modsatte, men ogsaa aldeles
modsigende.“

Dr. M., der til en vis Grad --- men ogsaa blot relativt
og til en vis Grad --- vil have „det Subjective, og
ligeledes til en vis Grad --- men ogsaa „blot relativt og
til en vis Grad“, vil have det Objective, faaer da hverken
den objective eller den subjective Afgjørelse, men en neutraliserende
Forening, der er --- principløs.

Hvad har nu Forfatteren af dogm. Oplysn. gjort for
at belyse dette Punkt? Har han nogetsteds forsøgt at tydeliggjøre
os, hvorledes en Erkjendelse, der er saa subjectiv,
saa inderlig, saa absolut personlig, som en sand Troeserkjendelse
nødvendigviis maa være, kan omsættes i Speculationens
og Videnskabens objective, upersonlige Form, og endda
være Troeserkjendelse? Nei, Hs. Høiærv. er kommen langt
nemmere til sin dogmatiske Oplysning, idet han (S. 42)
har gjort opmærksom paa, at det, jeg kalder den existentielle
Tænkning, vel omtrent maa være „det Samme, hvad han i
sin Dogmatik ogsaa har kaldt den existentielle Tænkning og
med Daub og Mynster kaldt Erkjendelsen i Religionen.“

Saaledes staaer da nu denne Sag. Man har tilladt
sig at recensere „den christelige Dogmatik“, hvis Forord
% --- p. 63 ---
\setcounter{footnote}{0}%
\opage{63}just i denne Henseende saae saa indbydende ud. Hvad de
tvende her omhandlede Recensioner (Mag. Sts og min)
angaaer: afviser Dr. Martensen den første, indlader sig ---
som viist er --- ikke paa den sidste, men erklærer reent ud,
at Kritiken over hans Dogmatik „endnu ikke egentlig kan
siges at være begyndt.“ Man har i disse Recensioner stillet
Dr. M. bestemte Theses, efterviist ham i hans Dogmatik
bestemte Modsigelser, anført ham bestemte Grunde. Dr. M.
erklærer det Hele for et tomt „Formalisterie“ og appellerer
til Stoffet.

Men, uagtet Dr. M., trods al sin Bestemthed, som
sagt, ikke vil gaae ind paa nogen bestemt Besvarelse af
noget bestemt Hovedspørgsmaal, vil han dog ikke destomindre
gjøre det „uomstødeligt vist, at man ikke har forstaaet ham,
men paa det  Vilkaarligste misforstaaet og mistydet hans
Tanker.“ --- Forhandlingen gaaer altsaa følgende Gang:

Man gjør Dr. M. opmærksom paa, at det i Troeserkjendelsen
bliver nødvendigt at holde sig absolut til Troens
Princip: Hs. Høiærv. forstaaer det, som om man vilde udgive
sig for „absolut i Troen.“

Man gjør Dr. M. opmærksom paa, at man i Videnskaben
absolut maa holde sig til Videns. Princip. Hs. Hv.
forstaaer det, som om man selv vilde udgive sig for „absolut
i Viden.“

Man gjør Dr. M. opmærksom paa, at der er en
Uendelighedens Forskjel imellem Troens Interesse og Videns
Interesse: Hs. Hv. forstaaer det, som om man talte
til ham om en „ligegyldig og interesseløs Viden.”

Man gjør Dr. M. opmærksom paa, hvorledes den
Speculation, „han har vedkjendt sig i sin Dogmatik“, oplø%
% --- p. 64 ---
\setcounter{footnote}{0}%
\opage{64}ser sig ved sin egen indre Modsigelse: Hs. Hv. forstaaer
det, som om man aldeles ikke tog Hensyn til „den Speculation“,
han har vedkjendt sig i sin Dogmatik.

Man gjør Dr. M. opmærksom paa, at et systematisk
Hoved (være sig Philosoph eller Theolog), der vil hævde
„den sammenhængende Erkjendelse“ og revse Dem, der maae
nøies med „Strøtanker, Aphorismer, Indfald og Glimt“,
fremfor Alt maa lægge sig efter at have et Princip og respectere
dets Conseqvenser: Hs. Hv. forstaaer det, som om
man herved begik en --- „\textit{metabasis in aliud genus}!“

Man gjør Dr. M. opmærksom paa, at han umulig
kan faae Speculationen med, hvis han virkelig indrømmer,
at Troens Tænkning er „en existentiel Tænkning.“ Hs. Hv.
forstaaer det, som om man havde glemt, at ogsaa han i sin
Dogmatik indrømmer, at Troens Tænkning er existentiel.

Man gjør Dr. M. opmærksom paa, at Troen ganske
vist har sin Gjenstand, sit bestemte Object, men at Sandhedserkjendelsen
i Troen netop forholder sig til dette Object
saaledes, at den subjective Vished, den subjective Tilegnelse
her bliver det ene Afgjørende: Hs. Hv. forstaaer det, som
om man talte til ham om en Sandhedserkjendelse, der
intet Objectivt havde, og belærer os (Oplysn. S. 16), at
„selv den fattigste rationalistiske Dogmatik“ dog maa erkjende,
„at Sandheden har en vis Objectivitet.“

Man gjør Dr. M. opmærksom paa, at efterdi Troens
Gjenstand nødvendigviis maa svare til Troen og Troeserkjendelsen,
saa kan Troens Gjenstand just af denne Grund
ikke være en Videns Gjenstand, da Troen som Tro netop
er det Modsatte af objectiv Viden: Hs. Hv. forstaaer det
% --- p. 65 ---
\setcounter{footnote}{0}%
\opage{65}(jfr, Oplysn. S. 19—22), som om man talte til ham om
en Tro, der skulde befinde sig i Strid med sin Gjenstand.

Man gjør endelig Dr. M. opmærksom paa, at Troesgjenstandenes
objective Ubegribelighed just bliver og maa
blive indlysende for den, der veed, hvad det vil sige, at begribe
Noget, og derhos alvorligt og vedholdende tænker over
sin Tro. Dr. M. forstaaer det, som om man nu var bleven
kjed af at tænke over sin Tro, og udbryder med Anselmus:
„\textit{negligentiæ mihi videtur},\dots{} Ja: \textit{negligentiæ mihi videtur}!“
---

Denne er da i Hovedsagen Dr. Ms. Forstaaelse af sine
„Modstanderes“ Tanker. Ja, hvad skal man sige? Skal en
saadan Forstaaelse nu ogsaa være et Tegn paa „Tankens
Forløsning“: da seer det --- idetmindste for mine Øine ---
noget broget ud med „Gjenfødelsens Logik.“

Om de Forsøg, Forf. af „dogm. Opl.“ her og der
har gjort paa at forklare og understøtte sit saa „mistydede“
Videnskabsbegreb vil jeg nu fatte mig i al Korthed,
a) Den „dogmatiske Speculation“ kan ikke forsvares
ved en Paaberaabelse paa „Aabenbaringsbevidstheden“, thi
Aabenbaringens Kjendsgjerninger ere --- som vi have seet
--- kun for Troens existentielle Tilegnelse: en speculativ
Aabenbaringsbevidsthed er en Modsigelse.

b) Den „dogmatiske Speculation“ kan ikke forsvares
ved Henviisning til den „Contemplation“, der findes i Davids
Psalmer og i Pauli Breve; at speculere er ikke det
Samme, som at synge Psalmer; at omvende Hedninger er
ikke det Samme som at speculere over Treenigheden: en
videnskabelig Psalmendigter, en speculerende Apostel, er en
Modsigelse

% --- p. 66 ---
\setcounter{footnote}{0}%
\opage{66}c) Den „dogmatiske Speculation“ kan ikke forsvares
ved Henviisning til de gamle Fædre, Tertullian, Augustin,
Anselm o. s.v.: de kjære Fædre kunde gjerne være gode
Christne, og endda befinde sig i nogen Uklarhed over Forholdet
imellem „\textit{certitudo experientiæ}” og „\textit{certitudo}
\textit{speculativa}“, ligerviis som ogsaa den kjære Dr. Martensen,
der endnu, trods Alt, hvad Aarhundredernes videnskabelige
Tænkning har oplyst, synes at ville forstene sig i samme
Uklarhed, dog desuagtet for sin Person gjerne kan være en
god Christen.

d) Den „dogmatiske Speculation“ kan endelig heller
ikke forsvares ved Henviisning til Troeslærdommenes indbyrdes
Sammenhæng. Kun speculative Sætninger ere i
speculativ Sammenhæng. Skulde Sammenhængen imellem
Dogmerne være speculativ, da maatte jo hvert enkelt Dogme
(altsaa ogsaa Dogmet om den forklarede Christi Allestedsnærværelse)
være speculativt; men hermed forholder det sig
nu, som vi have seet, anderledes: en speculativ Troesbekjendelse,
en speculativ Catechismus, en speculativ Dogmatik ere
Modsigelser, som end ikke „Gjenfødelsens Logik“ skal kunne
„mediere.“

Dog, som sagt, jeg fatter mig kort, deels fordi jeg nu
antager, at idetmindste Stridspunctet maa være synligt for
dem, der kunne og ville see, deels ogsaa fordi jeg ikke kan
slutte denne Forhandling uden endnu engang at henlede
Læserens Opmærksomhed paa Dr. Ms. og min Stilling til
„Pseudonymen“ og den Kierkegaardske Litteratur.

Istedetfor at holde sig til min Anmeldelse af „den christelige
Dogmatik“, har da Forf. af dogm. Opl. foretrukket
at holde sig til mit Uddrag d. e. til min Anmeldelse af
% --- p. 67 ---
\setcounter{footnote}{0}%
\opage{67}„Johannes Climacus.“ --- Ved saaledes at slaae tvende saa
ueensartede Størrelser, som „Pseudonymen“ og min Ringhed
sammen i Eet, har Hs. Hv. uden Tvivl tilsigtet at slaae,
som man siger, „tvende Fluer med eet Smæk,“ idet han
vistnok har meent at kunne slaae mig i „Pseudonymen“ og
saa med det Samme slaae „Pseudonymen“ i mig.

Har nu Dr. M. ramt mig i „Pseudonymen“? --- Jeg
kan ikke fornemme det. Thi hvad det angaaer, at jeg har
talt i min „Mesters Ord“, da kunde jeg jo ikke vel gjøre
Andet, naar jeg vilde give et paalideligt „Excerpt“ af hans
„Læresætninger.“ Hvad det angaaer, at jeg ikke engang har
søgt at forklare Hs. Hv., hvad jeg tænker mig ved „\emph{Paradoxet}“:
da bliver dette kun at forstaae saaledes, at Hs. Hv.
ikke engang har forsøgt at gjennemtænke, hvad der udtrykkelig
staaer i mine Skrifter. Hvad endelig det angaaer, at
jeg af „Letsindighed“ og „dilettantmæssige Ustadighed“ skulde
have stiftet Standpunct, blot for at høre blandt de „særdeles
Interessante“, da maatte det vel først godtgjøres, at mit
valgte „Standpunct“ er uholdbart. Det tilsigtede „Smæk“
kan altsaa ikke ramme mig, før det har ramt „Pseudonymen.“


Har da Dr. M. virkelig ramt „Pseudonymen“ i mig?
--- Jeg vil ikke fornærme Mag. Kierkegaard og gjøre mig
latterlig i „Pseudonymernes“ Øine ved at forsvare „Johannes
Climacus“ imod Dr. Martensen; jeg vil kun med et Par Ord
belyse Dr. Ms Mening om hvad han kalder Pseudonymernes
„vidtløftige Litteratur“ ved at henholde mig til de i hans
„dogmatiske Oplysninger“ henkastede Ytringer.

Alt i en Række af Aar have nu de pseudonyme Tænkere,
den Ene efter den Anden, fæstet Opmærksomheden paa
% --- p. 68 ---
\setcounter{footnote}{0}%
\opage{68}„Inderlighedens“ Opgaver og opbudt al deres Dialektik paa
at gjøre det christelige Problem kjendeligt; alt i en Række
af Aar har Dr. M., der som Dogmatiker beskjæftiger sig med
det samme Problem, været vidende om, at disse Pseudonymer,
skjøndt altfor indadvendte til at yttre sig med nogen
endelig Personlighedspolemik, dog alligevel bestandig holdt
et vaagent Øie med „Speculationen“ og „Systemet“; alt i
en Række af Aar har en stille voxende Opmærksomhed fulgt
denne saa forunderlig frastødende og dog saa uendelig tiltrækkende
„Litteratur,“ saa at Mange vistnok mange Gange maae
have spurgt sig selv: hvad er herved at gjøre? --- Under
disse Omstændigheder maae vi da være Dr. M. dobbelt Tak
skyldige, at han endelig træder op, og ikke blot siger os,
hvad herved er at gjøre, men tillige siger os det i en saa
„ligefrem Meddelelse.“ Af denne „ligefremme Meddelelse“
maa dog vel Enhver kunne see, at Dr. M. ikke har læst og
vist heller ikke agter at læse „Pseudonymernes Skrifter“, at
Hs. Hv. ikke kan forstaae disse Skrifter\footnote[1]{Dr. Martensens bestemte Erklæring, at han ikke kan forstaae „Pseudonymerne“, turde dog her muligviis være at tage i en mere bogstavelig Mening, end Hs. Hv. egentlig ønsker det. Hvad han saaledes f. Ex. (Opl. S. 39), vistnok ikke saameget efter Anselms som efter Frederikke Bremers „Forbillede“, forklarer os om Støttehelten, „den Enkelte,“ item de sindrige Betragtninger, han her og der anstiller over Troens „Lidenskab“ og Troens „Risico“, synes virkelig at tyde hen paa, at hans „dogmatiske Genius“ i denne Henseende er lovlig undskyldt. Hvad der her billigviis taler til Dr. Ms. Undskyldning, er da nu ogsaa dette, at Pseudonymerne aldeles ikke tilbyde ham noget fast og tilforladeligt ποῦ στῶ.\par De dialektiske Pseudonymer have, som det synes, alle indrettet sig efter „Stadiernes“ Lichtenbergske Motto: „Solche Werke sind Spiegel; wenn ein Affe hinein guckt, kann kein Apostel heraus sehen.“ Et saadant „Standpunct“ er unegtelig kun lidet „dogmatisk“; men --- man faaer at tage det, som det er, da andet ποῦ στῶ ikke tilbydes.}, „at han ifølge sine
% --- p. 69 ---
\setcounter{footnote}{0}%
\opage{69}Studiers Gang“ og sin „individuelle Aandsretning“ just
heller ikke ønsker at forstaae disse Skrifter, og at altsaa det
Hele heller ikke kan have stort at betyde.

Det er et haardt Slag, ja, det er et dræbende Slag,
det Slag, hvormed Dr. Ms. dogmatiske Overlegenhed her
har slaaet den kun lidet dogmatiske „Johannes Climacus“;
men --- det maa jeg tilstaae --- uagtet dette Slag, hvad
Realiteten angaaer, er et saa haardt Slag, har Dr. M., hvad
Formen angaaer, alligevel vidst at „iagttage Grændsen“ og
føre Slaget med al den Sindighed og tilsyneladende Skaansomhed,
der altid klæder den Overlegne.

Dr. M. erklærer jo rigtignok, at han aldeles ingen Brug
har for de „Læresætninger“, jeg af „Pseudonymen“ har villet
„meddele“ ham, efterdi han paa disse maa anvende det gamle
Ord, „at det Sande deri ikke er Nyt, og det Nye ikke Sandt“;
men han indrømmer dog tillige, at „Pseudonymen“ kan til
en vis Grad være at undskylde dermed, at han som „Gymnastiker“
og „Humorist“ blot har til Hensigt „paa sokratisk
Viis at sætte en dybere Skepticisme i Bevægelse (Oplysn.
S. 12).“ Den meste Skyld falder altsaa paa mig, der har
været saa „letsindig“ at tage Humoristens „Sætninger“ for
ramme Alvor, og mene, at disse Sætninger kunde (Oplysn.
S. 13) bruges til „en ny Erkjendelseslære og en derpaa
begrundet Behandling af Theologien.“ --- Men, som sagt,
denne Skaansomhed mod „Gymnastikeren“ ligger alene i
Formen. Den dogmatiske Dom lyder saaledes:

Skal Det, „Johannes Climacus“ siger om det christelige
Problem, forstaaes som en gymnastisk-humoristisk Halvspøg,
kan der maaskee gjerne være Noget i det; men dersom hans
Hovedsætninger skulle tages for Alvor og gjælde for Sand%
% --- p. 70 ---
\setcounter{footnote}{0}%
\opage{70}hed: da maa Dr. M. paa det Bestemteste erklære, at „det
Sande den ikke er Nyt, og det Nye ikke Sandt“; og altsaa
gaaer det jo dog i Realiteten alligevel ud over Pseudonymen.


Hvad nu min Mening angaaer, siger jeg til alt Dette
ganske ligefrem: læs de Kierkegaardske Skrifter ! Læser, Du,
som vil have nogen Dom i denne Sag, hvo Du end er,
lærd eller ulærd, læs blot de Kierkegaardske Skrifter i al
Stilhed uden at forstyrres af Dogmets Surren og Partiernes
Strid. Læs disse Skrifter med intellectuel Interesse!
prøv hvad Du læser, ved Dig selv alene, og Du skal see,
hvor let og spøgende, og dog --- hvor alvorlig, de „humoristiske“,
Forfattere forene sig om at føre Dig ind i „Problemet“, i
Inderlighedens, i Personlighedens, i Evighedens sande Problem,
i det Problem, for hvis Skyld det dog alene er værd
at leve. Eller: læs disse Skrifter med religiøs-christelig
Interesse! frigjør dit Sind fra alle Endelighedens udadvendte,
verdslig-kirkelige Lidenskaber, anvend det Læste paa
Dig selv, og Du skal see, med hvilken instinctagtig og dog bevidst
Sikkerhed de ellers saa forskjellige Forfattere forstaae at vise
Dig det Christelige i Christendommen, og hvilken Mildhed Forsoning
og Fred, der dog her oversvæver de mange ellers saa
ængstendende og dybt forfærdende Tanker. Eller, hvis Du
hellere vil det saaledes: læs da disse Skrifter i afgjort polemisk
Interesse! Antag, at hvad jeg her har sagt, blot er en
taabelig Overdrivelse, udsprungen enten af blind Enthusiasme
eller underfundig Smigren for et andet Menneske, hvis Anerkjendelse
maaskee til Gjengjæld forventes, om end i ringere
Maal; antag, at det Hele ikke er andet end en vilkaarlig
trodsende „Eensidighed,“ der paa det Eftertrykkeligste maa
% --- p. 71 ---
\setcounter{footnote}{0}%
\opage{71}gjendrives; sæt Dig alvorlig for, at Du netop vil være den, der
møder med Gjendrivelsen, vælg dine bestemte Puncter, gjør
et Udkast, men --- læs imidlertid alle de paagjældende Skrifter,
„Pseudonymerne“ ligesaavel som de „opbyggelige Taler“, læs
dem grundig (thi, hvis Du ikke læser grundig, kan Du vist
heller ikke gjendrive grundig): Du skal see, at det endda
just ikke gaaer saa ganske let med Gjendrivelsen; det turde
let hændes Dig, at Du ved at gaae ind i den virkelige
Sammenhæng her forefandt en Selvkritik, der gjorde dine
Indsigelser overflødige, og tilsidst lagde din hele paatænkte
Kritik under Beslag.

At jeg forøvrigt ikke er i Besiddelse af nogen ligefrem
eller authentisk Forklaring af disse Skrifter, er en Selvfølge,
da enhver ligefrem Forstaaelse her er gjort dialektisk umulig.
At Umuligheden af en saadan ligefrem Forstaaelse
skulde tjene som Beviis for, at den „indirecte Meddelelse“
er ufuldkommen og af lavere Art end den directe, er imidlertid
en feilagtig Mening, der er let at corrigere.

Man har jo rigtignok meent, at „Pseudonymerne“ med
deres humoristisk spøgende Form og deres dobbelt reflecterede
Kategorier skulde være mindre alvorlige og derfor mindre
tilforladelige, end de „opbyggelige Taler“ (f. Ex. „Kjærlighedens
Gjerninger„). Det Modsatte er just Tilfældet. Netop ved % PRINTED AS IS: the closing quote printed low
deres indirecte Meddelelsesform staae disse „Pseudonymer“
just i Alvor paa lige Trin med de opbyggelige Skrifter\footnote[1]{Hvis Meddelelse desuden, skjøndt i en ganske anden Forstand, er og maa være indirect.},
(ja, ere med al deres Vid og Humor i Grunden langt strengere);
ved deres „indirecte Meddelelsesform“ opnaae Pseu%
% --- p. 72 ---
\setcounter{footnote}{0}%
\opage{72}donymerne netop, hvad der er aldeles nødvendigt, naar der
skal tales absolut om det Absolute, de opnaae nemlig ved
deres Forvandlinger i det Komiske, at afskjære al relativ
Raisonneren og standse den ellers ustandselige Discussion.


Havde f. Ex. „Johannes Climacus“ foredraget sine af mig
anmeldte „Læresætninger“ i en reen og ligefrem videnskabelig
Form, det vilde være blevet Theologerne en alvorlig
Sag, og de vilde da sikkerlig ogsaa have taget det meget
alvorlig. Nu derimod, da han har meddeelt disse Læresætninger
i Vittighedens og Humorets indirecte Form, nu anseer
man ikke Sagen for saa alvorlig. Heri sees en Forvexling
af det Høiere og det Lavere.

Med en blot videnskabelig Form vilde „Forfatteren“
aabenbart have bevæget sig i en Cirkel, forsaavidt han,
skjøndt polemiserende mod Videnskaben, dog alligevel forblev
indenfor Videnskaben, og saaledes forsøgte sig i den Modsigelse
at gjøre det Ikke-Videnskabelige videnskabeligt. Med
en blot videnskabelig Form vilde „Forfatteren“ have befundet
sig i den Modsigelse, at tale ud af Subjectivitetens og
Inderlighedens Grund uden dog at tale Subjectivitetens og
Inderlighedens Sprog, at sige det „ligefrem“, som dog intet
Menneske kan see eller forstaae „ligefrem.“ Naar da en Forfatter.
der vil sætte et Princip igjennem og gjøre et Problem
kjendeligt, netop bruger den Meddelelsesform, som ifølge
Sagens Natur, er den eneste for hans Problem adæqvate
skulle vi saa ikke være enige om at indrømme, at Valget
af en saadan Form just tyder paa, at denne Forfattere har % PRINTED AS IS: Forfattere (for Forfatter)
--- Alvor; og, naar det endelig indsees, at Inderlighedens
Problem kun lader sig meddele i Inderlighedens Form, og
% --- p. 73 ---
\setcounter{footnote}{0}%
\opage{73}at denne Form, hvor det gjælder en Position ligeoverfor
„Videnskaben“, nødvendigviis maa blive en Form, der, idet
den udfolder Subjectivitetens Indhold, tillige sikkrer dette
Indhold derved, at enhver Yttring humoristisk ophæver sin
umiddelbare objective Gyldighed, før desto stærkere at gjenlyde % PRINTED AS IS: før for for (p. 73; stroke checked at the image)
i Inderligheden selv og hvile i Subjectiviteten: skulle vi
saa ikke være enige om at indrømme, at der udkræves langt
rigere Aandsgaver og en langt høiere Productivitet før at % PRINTED AS IS: før for for (p. 73)
meddele Inderlighedens Problem i Inderlighedens Sprog,
end der behøves for at fremsætte sine Tanker i en ligefrem
videnskabelig Form.

Man vil af det her Anførte let kunne see, at jeg betragter
den Kierkegaardske Litteratur som noget i sig saa
Eiendommeligt, Afsluttet og Enestaaende, at her aldeles ikke
kan være Tanke om, at faae disse Skrifter ad nogen ligefrem
Vei indførte og optagne med i den øvrige Litteratur. Men
deraf følger ingenlunde, at man i „Litteraturen“ skulde have
Lov til at ignorere disse Skrifter eller lade, som om man
ikke havde læst dem. Tvertimod! Enhver, der nuomstunder
vil tale med om Forholdet mellem Troen og Erkjendelsen
(han være Theolog eller Philosoph, Præst eller Læg), skal
visselig ikke opklare os Meget, hvis han selv befinder sig
i Uklarhed over disse Skrifters „Problem.“

Forsaavidt de Kierkegaardske Skrifter tilhøre Offentligheden,
tilhøre de naturligviis ogsaa Litteraturen. Enhver
Forfatter, hvis Opgave berører deres, kan og bør sætte sig
i et aldeles bestemt og utvetydigt Forhold til dem.

Om dette Forhold skal være polemisk afvisende eller velvillig
anerkjendende, derom spørge Enhver sig selv, uden at
tage Mag. Kierkegaard med paa Raad. Har Mag. Kierke%
% --- p. 74 ---
\setcounter{footnote}{0}%
\opage{74}gaard saa Noget at indvende, kan han jo sige til, og selv
gaae med i Discussionen.

At jeg for min Deel ikke blot har personlig Velvillie for
Mag. Kierkegaard men ogsaa stor Høiagtelse for den Mand
af hvis Skrifter jeg tydelig nok (om end ikke „ligefrem“)
kan see, at hans indre Liv maa være langt mere anstrenget
end mit, behøver jeg vel ikke yderligere at forsikkre. Men
det skulde dog gjøre mig inderlig ondt, om Mag. Kierkegaard
(hvad jeg vel maa ansee for utænkeligt) skulde misforstaae
disse mine Yttringer, som om jeg herved vilde sætte
mig i nogetsomhelst Afhængighedsforhold til hans personlige Individualitet. Mit Forhold til denne Forfatterindividualitet
er og maa væsentlig være det samme, som det
vilde være, hvis han, istedetfor at leve her i Kjøbenhavn,
levede t. Ex. i Rom eller i Philadelphia, eller som, hvis han, % PRINTED AS IS: t. Ex. (for f. Ex.)
istedetfor at være vor „høistærede Samtidige“, antoges at have
levet enten med Tertullian eller med Augustinus eller med
--- Anselmus.

At komme i et usandt Forhold til de Kierkegaardske
Skrifter anseer jeg for yderst farligt, forsaavidt jeg altid
anseer det for yderst farligt at komme i et usandt Forhold,
til en Sandhed; at komme i en litterær Polemik med Mag.
Kierkegaard anseer jeg derimod i denne Forstand ikke for
saa farligt. Lad ham have nok saamegen Kløgt og nok saa
glimrende Vaaben; Herregud, vi kunne jo ikke alle være Genier;
naar man strider for en Sag, strider man jo ikke for
at vise sit Talent: Mag. Kierkegaard er desuden, saavidt jeg
veed, netop den eneste Forfatter, der har opbudt, ja jeg kunde
sige: offret og ødslet alt sit Talent paa at bevise, at for
Erkjendelsen af den høieste Sandhed er Talentet et Lige%
% --- p. 75 ---
\setcounter{footnote}{0}%
\opage{75}gyldigt. Lad Mag. Kierkegaard paavise, at der er enten
Tvetydighed, Uklarhed eller anden Mislighed i den Brug,
jeg har gjort af hans Skrifter: kan han overbevise mig,
velan, saa tilstaaer jeg mit Feilgreb og takker ham for hans
Veiledning; kan han ikke overbevise mig, nu, saa bliver jeg
ved Mit, og skulde mene, at jeg i saa Henseende har samme
Ret, som enhver Anden.

Altsaa: min Sag er ikke Mag. Kierkegaards Sag.
Mag. Kierkegaards Sag er udenfor al Discussion, ideelt
befæstet i det Inderlighedens Middelpunkt, hvorom det Religiøse
bevæger sig. Min Sag er i Discussionen, udadvendt
mod det Grændseskjel, hvor Theologie og Philosophie
strides om Viden og Ikke-Viden.

Da jeg udgav mit første Skrift: „Evangelitroen og % PRINTED AS IS: Evangelitroen
den moderne Bevidsthed“, meente jeg at kunne føre min
Sag imod Theologien, uden at berøre Nogen af vore Theologer
personlig. Dr. Martensens Holdning i Forordet til
hans Dogmatik opfordrede mig til at anmelde denne Bog
i en saadan Tone, at hans fornemme Tvetydighed nødvendigviis
maatte briste. Imidlertid har den ene Theolog følt
sig fornærmet efter den Anden; men, skjøndt jeg i mit
andet Skrift: „Evangelietroen og Theologien,“ har gjort udførligt
Rede for min Opfattelse af Christendommens Forhold
til Videnskaben, har dog ingen Theolog endnu forsøgt enten
at omstøde mine „Sætninger“ eller overhovedet at sætte sig
„videnskabelig“ ind i Sagen.

\end{document}
